
\chapter{روش پیشنهادی}
\label{chap:method}

این فصل روش پیشنهادی این پایان‌نامه را ارائه می‌کند. ساختار فصل مستقیماً از چهار جای خالی‌ای پیروی می‌کند که در پایان فصل~\ref{chap:literature} استخراج شد: نبود سازوکار قابل‌اعتماد برای تشخیص فاز کنتراست مستقل از سرآیند، نبود پایپ‌لاین ثبت تصویری که \emph{دامنه‌ی} محاسبه‌ی معیار تراز را به نواحی پایدار محدود کند، کمیابی جفت‌های هم‌تراز برای آموزش مدل‌های سنتز، و فقدان شواهد کمّی درباره‌ی اثر این مداخلات بر شکاف دقت بخش‌بندی.

بخش~\ref{sec:m-overview} نمای کلی و منطق زنجیره‌ی مؤلفه‌ها را شرح می‌دهد. بخش~\ref{sec:m-data} داده و پیش‌پردازش مشترک را معرفی می‌کند. سه بخش بعدی به‌ترتیب سه مؤلفه‌ی روش را ارائه می‌دهند: تشخیص فاز کنتراست (بخش~\ref{sec:m-phase})، ثبت تصویر بین‌فازی چندگذره (بخش~\ref{sec:m-reg}) و سنتز فاز کنتراست‌دار (بخش~\ref{sec:m-syn}). بخش~\ref{sec:m-eval} پروتکل ارزیابی را تعریف می‌کند؛ این بخش صرفاً فهرستی از معیارها نیست، بلکه خود بخشی از سهم روش‌شناختی این پایان‌نامه است، چرا که—چنان‌که نشان داده خواهد شد—معیارهای متعارف سطح‌پیکسل قادر به رتبه‌بندی مدل‌های این مسئله نیستند.

\section{نمای کلی}
\label{sec:m-overview}

سه مؤلفه‌ی این روش یک زنجیره‌ی وابسته تشکیل می‌دهند و نه سه کار موازی.

نقطه‌ی شروع، محدودیت داده است. برای آموزش نظارت‌شده‌ی هر مدل سنتز \lr{NC}-به-\lr{CE}، به جفت‌هایی از یک بیمار نیاز است که هم فاز بدون‌کنتراست و هم فاز کنتراست‌دار آن موجود و از نظر هندسی هم‌تراز باشند. مرور مجموعه‌داده‌ها~\cite{danaei2026abdomenctreview} نشان داد چنین داده‌ای کمیاب است. بنابراین جفت‌ها باید از داده‌ی چندفازی موجود \emph{ساخته} شوند، و ساختن آن‌ها دو پیش‌نیاز دارد:

\begin{itemize}
\item \مهم{شناسایی فاز.} برای آن‌که بدانیم کدام سری از یک بررسی، فاز بدون‌کنتراست و کدام فاز وریدی است، به برچسب فاز نیاز است. سرآیند \lr{DICOM} این برچسب را به‌طور قابل‌اتکا فراهم نمی‌کند (بخش~\ref{sec:phase-detection}). مؤلفه‌ی نخست این خلأ را پر می‌کند.
\item \مهم{هم‌ترازسازی هندسی.} دو سری از یک بررسی در دو زمان و با تنفس متفاوت گرفته شده‌اند و بدون ثبت تصویر، جفت \emph{نیستند}. مؤلفه‌ی دوم این هم‌ترازی را انجام می‌دهد.
\end{itemize}

مؤلفه‌ی سوم، سنتز، بر خروجی این دو تکیه می‌کند. اما رابطه‌ی مؤلفه‌ی نخست با مؤلفه‌ی سوم عمیق‌تر از یک پیش‌نیاز اجرایی است: مدل تشخیص فاز، علاوه بر نقش دروازه‌بانی داده، به \emph{ابزار اندازه‌گیری} ارزیابی سنتز نیز تبدیل می‌شود. همان‌طور که در بخش~\ref{sec:m-eval} تشریح می‌شود، معیار اصلی این پایان‌نامه برای سنجش وفاداری کنتراستِ تصویر تولیدشده، از ویژگی‌های سطح‌اندام همین مدل استخراج می‌شود. به این ترتیب مؤلفه‌ای که برای آماده‌سازی داده ساخته شد، مسئله‌ی سنجش را نیز حل می‌کند.

سرانجام، پرسش پایانی پایان‌نامه—آیا این مداخلات شکاف دقت بخش‌بندی چندعضوی را کاهش می‌دهند—یک ارزیابی پایین‌دستی\footnote{\lr{Downstream}} جداگانه می‌طلبد که در بخش~\ref{sec:m-downstream} تعریف می‌شود.

% شکل ۳-۱: نمودار جریان کل روش. سه بلوک (تشخیص فاز / ثبت تصویر / سنتز)
% به‌علاوه‌ی بلوک ارزیابی پایین‌دستی، با فلش داده میان آن‌ها. باید نشان دهد
% که مدل تشخیص فاز دو بار ظاهر می‌شود: یک‌بار به‌عنوان دروازه‌ی داده و
% یک‌بار به‌عنوان ابزار اندازه‌گیری.

\section{مرور نظام‌مند مجموعه‌داده‌های سی‌تی شکمی: پروتکل}
\label{sec:m-review}

بخش~\ref{sec:datasets} نشان داد هیچ ارزیابی نظام‌مندی از مجموعه‌داده‌های عمومی سی‌تی شکمی به‌عنوان منبع آموزش هوش مصنوعی وجود ندارد. این بخش، پروتکلی را که برای پرکردن این شکاف تدوین شد شرح می‌دهد؛ یافته‌های حاصل در بخش~\ref{sec:e-review} گزارش می‌شوند. نتایج این مرور به‌طور مستقل نیز منتشر شده است~\cite{danaei2026abdomenctreview}.

\subsection*{شناسایی و معیارهای ورود}

جست‌وجو در \lr{Google Scholar}، \lr{PubMed}، \lr{Scopus} و مخازن نهادی از جمله \lr{TCIA} انجام شد. جست‌وجوی نخست میان مهر و آذر \عدد{۱۴۰۳} و دو جست‌وجوی به‌روزرسانی میان آذر \عدد{۱۴۰۴} و بهمن \عدد{۱۴۰۴} صورت گرفت. افزون بر جست‌وجوی پایگاهی، فهرست مراجع و توصیف منبع داده‌ی هر مجموعه‌داده‌ی یافت‌شده بررسی شد تا مجموعه‌داده‌های سازنده‌ای که در مجموعه‌های ترکیبی بازاستفاده شده‌اند ردیابی شوند---گامی که برای کمّی‌کردن بازاستفاده ضروری بود. مخازن خارج از ادبیات متعارف (\lr{Kaggle}، \lr{Zenodo}، \lr{Grand Challenge}) نیز پیمایش شدند.

پنج معیار ورود اعمال شد: داده باید حجمی\footnote{\lr{Volumetric}} سه‌بعدی باشد (نه برش‌های دوبعدی استخراج‌شده)؛ دست‌کم یک ناحیه‌ی برچسب‌خورده به‌صورت نقاب چندضلعی یا وکسلی داشته باشد؛ برچسب‌ها از وظیفه‌ای بالینی پشتیبانی کنند؛ برچسب‌ها به ساختارهای شکمی مربوط باشند؛ و مجموعه‌داده دست‌کم در یک مقاله‌ی داوری‌شده استفاده شده باشد. معیار آخر بدین معناست که این مرور، \emph{چشم‌انداز منتشرشده} را نمایندگی می‌کند، نه کل منابع موجود.

جست‌وجو \عدد{۴۸} نامزد شناسایی کرد که پس از اعمال معیارها سه مورد کنار گذاشته شد و \textbf{\عدد{۴۵} مجموعه‌داده} باقی ماند.

\subsection*{استخراج داده}

دو ارزیاب مستقل، داده را با فرم استاندارد استخراج کردند و در موارد اختلاف، ارزیاب سوم داوری کرد. استخراج در شش محور انجام شد: ترکیب مجموعه‌داده، بافت بالینی، خاستگاه داده، ویژگی‌های تصویربرداری (شامل \textbf{وضعیت فاز کنتراست})، جزئیات برچسب‌گذاری، و اطلاعات جمعیت‌شناختی. محور چهارم، که وضعیت فاز کنتراست را ثبت می‌کند، همان محوری است که ارقام مورد نیاز بخش‌های~\ref{sec:m-reg} و~\ref{sec:m-syn} این فصل را فراهم می‌آورد.

\subsection*{چارچوب ارزیابی ریسک سوگیری}

از آن‌جا که هیچ چک‌لیست پذیرفته‌شده‌ای برای این منظور وجود ندارد (بخش~\ref{sec:datasets})، یک چارچوب هفت‌مقوله‌ای با تطبیق سه ابزار موجود تدوین شد: \lr{PROBAST+AI}~\cite{Moonse082505}، معیارهای کیفیت کوچاک و همکاران~\cite{Kocak2024}، و راهنمای گزارش‌دهی \lr{TRIPOD+AI}~\cite{Collins2024}. هر مقوله در سه سطح ریسک پایین، متوسط یا بالا امتیاز می‌گیرد، و مقوله‌هایی که مستندات اولیه اطلاعات کافی برای آن‌ها فراهم نمی‌کرد با برچسب «در دسترس نیست» ثبت شدند. جدول~\ref{tab:bias-categories} این هفت مقوله را خلاصه می‌کند.

این ارزیابی تنها روی \عدد{۲۲} مجموعه‌داده با صد نمونه‌ی عمومی یا بیشتر انجام شد؛ مجموعه‌داده‌های کوچک‌تر به‌طور پیش‌فرض در اغلب مقوله‌ها ریسک بالا می‌گیرند و این امر ارزش تمایزی طبقه‌بندی را از میان می‌برد.

\begin{table}[H]
\centering
\footnotesize
\caption{هفت مقوله‌ی چارچوب ارزیابی ریسک سوگیری، به همراه تعریف عملیاتی و تهدید مرتبه‌بالاتری که هر مقوله به آن مربوط است.}
\label{tab:bias-categories}
\begin{tabular}{|p{2.6cm}|p{6.6cm}|p{3.6cm}|}
\hline
\textbf{مقوله} & \textbf{محور ارزیابی} & \textbf{تهدید مرتبه‌بالاتر} \\
\hline
سوگیری طیف\footnote{\lr{Spectrum bias}} & پوشش دامنه‌ی شدت و تنوع بالینی درون گروه مثبت & نمایندگی موارد \\
\hline
سوگیری انتخاب\footnote{\lr{Selection bias}} & پوشش تشخیص‌های افتراقی در گروه غیرمثبت & نمایندگی موارد \\
\hline
سوگیری برچسب‌گذاری\footnote{\lr{Labelling bias}} & دقت و سازگاری فرایند حاشیه‌نویسی؛ تعداد و تخصص بازبین‌ها & نمایندگی موارد و جابه‌جایی حوزه \\
\hline
سوگیری نژادی\footnote{\lr{Racial bias}} & نمایندگی گروه‌های نژادی و قومی نسبت به جامعه‌ی هدف & سوگیری جمعیتی \\
\hline
سوگیری جمعیت‌شناختی\footnote{\lr{Demographic bias}} & توازن جنسیت و توزیع سنی & سوگیری جمعیتی \\
\hline
سوگیری فنی\footnote{\lr{Technical bias}} & تنوع سخت‌افزار تصویربرداری، پروتکل و تزریق ماده‌ی حاجب & جابه‌جایی حوزه \\
\hline
سوگیری زمانی\footnote{\lr{Temporal bias}} & بازه‌ی زمانی گردآوری داده & جابه‌جایی حوزه \\
\hline
\end{tabular}
\end{table}

\subsection*{فرایند ارزیابی و پایایی}

سه ارزیاب---یک پژوهشگر تصویربرداری پزشکی، یک پزشک با تجربه‌ی پژوهشی در تحلیل تصویر پزشکی، و یک پزشک عمومی با تجربه‌ی پژوهشی در تصویربرداری---به‌طور مستقل همه‌ی مجموعه‌داده‌های واجد شرایط را ارزیابی کردند. پیش از ارزیابی مستقل، جلسه‌ی کالیبراسیون مشترکی برای یکسان‌سازی تفسیر تعاریف عملیاتی برگزار شد. ارزیابی تنها بر مبنای مستندات اولیه‌ی در دسترس عموم انجام شد و با سازندگان مجموعه‌داده‌ها تماس گرفته نشد؛ این تصمیم تضمین می‌کند ارزیابی، همان اطلاعاتی را بازتاب دهد که در اختیار جامعه‌ی پژوهشی است.

اختلاف میان دو یا چند ارزیاب، به‌عنوان تعارض تعریف و توسط ارزیاب چهارم در مشورت با یک رادیولوژیست بورد‌شده حل شد. پایایی بین‌ارزیاب با آلفای کریپندورف (با وزن‌دهی فاصله‌ی ترتیبی) سنجیده شد: مقدار کلی \عدد{۰٫۷۲۷} با فاصله‌ی اطمینان \عدد{۹۵} درصدِ \عدد{۰٫۶۲} تا \عدد{۰٫۸۴}، که توافق قابل‌ملاحظه تلقی می‌شود. مقادیر به‌تفکیک مقوله از \عدد{۰٫۱۹۹} برای سوگیری زمانی تا \عدد{۰٫۸۲۰} برای سوگیری انتخاب متغیر بود؛ پایین‌بودن توافق در سوگیری زمانی، بازتاب ناقص‌بودن گزارش بازه‌ی گردآوری در مستندات اغلب مجموعه‌داده‌هاست.

\subsection*{طبقه‌بندی کلی سوگیری}

برای تجمیع امتیازهای مقوله‌ای به یک طبقه‌بندی کلی، یک امتیاز پیوسته بر پایه‌ی سه اصل تعریف شد: ساختار چنددامنه‌ای \lr{ROBINS-I}~\cite{Sterne2016ROBINSI}، منطق کاهش تجمعی \lr{GRADE}~\cite{Schunemann2019GRADE}، و این اصل از \lr{QUADAS-2}~\cite{Whiting2011QUADAS2} که دامنه‌ی مستندنشده را نمی‌توان کم‌ریسک فرض کرد. مقوله‌های «در دسترس نیست» برای محاسبه، ریسک متوسط در نظر گرفته شدند:

\begin{equation}
\label{eq:bias-score}
\text{امتیاز سوگیری} = \sum(\text{ریسک بالا}) \times 1.0 \;+\; \sum(\text{ریسک متوسط} + \text{در دسترس نیست}) \times \tfrac{1}{3}
\end{equation}

آستانه‌ها متناسب با چارچوب هفت‌مقوله‌ای تعیین شدند: امتیاز \عدد{۳٫۰} و بالاتر «سوگیری بحرانی»، \عدد{۲٫۰} تا زیر \عدد{۳٫۰} «سوگیری بالا»، \عدد{۱٫۰} تا زیر \عدد{۲٫۰} «سوگیری متوسط»، و زیر \عدد{۱٫۰} «سوگیری پایین».

\section{داده و پیش‌پردازش}
\label{sec:m-data}

\subsection*{مجموعه‌داده}

آزمایش‌های این پایان‌نامه بر مجموعه‌داده‌ی عمومی \lr{VinDr-Multiphase}~\cite{vindr} انجام می‌شود. این مجموعه، بررسی‌های سی‌تی شکمی چندفازی را همراه با برچسب فاز در سطح سری فراهم می‌کند و از این رو هم برای آموزش و ارزیابی مدل تشخیص فاز و هم—پس از ثبت تصویر—برای ساخت جفت‌های آموزشی سنتز مناسب است.

انتخاب یک مجموعه‌داده‌ی عمومی به‌جای داده‌ی خام گردآوری‌شده، تصمیمی آگاهانه است: هر ادعای این پایان‌نامه درباره‌ی سنتز باید روی داده‌ای قابل بازتولید توسط دیگران آزمون شود، و مرور مجموعه‌داده‌ها~\cite{danaei2026abdomenctreview} نشان داد که همین قابلیت بازاستفاده، کمبود اصلی ادبیات این حوزه است.

پس از اعمال پایپ‌لاین ثبت تصویر (بخش~\ref{sec:m-reg}) و کنارگذاری جفت‌هایی که معیار کیفیت هم‌ترازی را برآورده نمی‌کنند، \lr{137} جفت هم‌ترازشده‌ی بدون‌کنتراست/وریدی و \lr{136} جفت بدون‌کنتراست/شریانی باقی می‌ماند. این جفت‌ها در سطح \emph{بیمار}—و نه در سطح برش یا سری—به سه بخش تقسیم می‌شوند: \lr{97} مورد آموزش، \lr{20} مورد اعتبارسنجی و \lr{20} مورد آزمون. تقسیم با بذر تصادفی ثابت انجام و در یک فایل منجمد ذخیره می‌شود، و \emph{همه‌ی} مدل‌های این پایان‌نامه—اعم از مدل‌های پیشنهادی و مدل‌های مبنای بازپیاده‌سازی‌شده از ادبیات—بدون استثنا روی همین تقسیم آموزش و ارزیابی می‌شوند. تقسیم در سطح بیمار، نشت اطلاعات میان برش‌های مجاور یک بیمار را که در ادبیات سنتز تصویر پزشکی خطای رایجی است حذف می‌کند.

\subsection*{استانداردسازی هندسی و شدت روشنایی}

همه‌ی حجم‌ها به فاصله‌ی وکسل همسان‌گرد\footnote{\lr{Isotropic}} \lr{1.5} میلی‌متر بازنمونه‌برداری می‌شوند. یکنواختی این مقدار در سراسر مجموعه به‌صورت تجربی تأیید شده است و نه فرض‌شده؛ اهمیت این تأیید در فصل~\ref{chap:experiments} روشن می‌شود، آن‌جا که ناهمگنی فاصله‌ی وکسل به‌عنوان یکی از توضیح‌های ممکن برای کف خطای سنتز کنار گذاشته می‌شود.

شدت روشنایی با یک پنجره‌ی \lr{HU} ثابت در بازه‌ی $[-200, 400]$ به بازه‌ی $[0,1]$ نگاشت می‌شود. این پنجره بافت نرم شکمی و محدوده‌ی کنتراست عروقی را پوشش می‌دهد. انتخاب یک پنجره‌ی \emph{واحد} برای همه‌ی مدل‌ها یک الزام روش‌شناختی است، نه یک جزئیات پیاده‌سازی: تغییر پنجره، مقدار عددی هر معیار سطح‌پیکسل را تغییر می‌دهد، و مقایسه‌ی دو مدل که با دو پنجره‌ی متفاوت سنجیده شده‌اند بی‌معناست.

\subsection*{نقاب‌های اندام}

نقاب‌های اندام با \lr{TotalSegmentator}~\cite{wasserthal2023} تولید می‌شوند. خروجی پرشمار این ابزار سپس به یک حجم برچسب فشرده با قرارداد ثابت ادغام می‌شود که در آن هر اندام یک شناسه‌ی عددی یکتا دارد. دو نقاب از این ادغام حاصل می‌شود و تمایز میان آن‌ها در ادامه‌ی فصل نقش محوری دارد:

\begin{itemize}
\item \مهم{نقاب کامل}: همه‌ی اندام‌های بخش‌بندی‌شده. این نقاب دامنه‌ی \emph{ارزیابی} است.
\item \مهم{نقاب ثبت}: تنها استخوان‌ها و اندام‌های نرمِ پایدار در برابر فاز کنتراست. روده، ساختارهای پرکنتراست عروقی و ساختارهای دارای حرکت دودی\footnote{\lr{Peristalsis}} به‌صراحت از آن حذف شده‌اند. این نقاب دامنه‌ی \emph{محاسبه‌ی معیار تراز} است.
\end{itemize}

نقاب‌ها از فاز کنتراست‌دار استخراج می‌شوند، چرا که مرزهای اندام در آن فاز به‌مراتب واضح‌تر است. پیامد این انتخاب آن است که این نقاب‌ها را نمی‌توان به‌عنوان ورودی به مولد سنتز داد—چنین کاری اطلاعات هدف را به ورودی نشت می‌دهد—و کاربردشان به وزن‌دهی تابع زیان و تعریف معیارها محدود می‌ماند. این محدودیت در بخش~\ref{sec:m-eval} دوباره مطرح می‌شود.

\section{تشخیص فاز کنتراست}
\label{sec:m-phase}

\subsection*{انگیزه‌ی طراحی}

چنان‌که در بخش~\ref{sec:phase-detection} مرور شد، بیشتر روش‌های موجود تشخیص فاز یا به شبکه‌های عصبی سه‌بعدی با نیاز داده‌ی آموزشی زیاد متکی‌اند یا مستقیماً روی حجم خام کار می‌کنند و در نتیجه به دستگاه، پروتکل و هسته‌ی بازسازی حساس‌اند. طراحی این پایان‌نامه از فرض متفاوتی آغاز می‌شود:

\begin{quote}
فاز کنتراست، در تعریف بالینی خود، چیزی جز \emph{الگوی شدت روشنایی در اندام‌های مشخص} نیست. رادیولوژیست فاز را از روی روشنایی آئورت نسبت به ورید باب و پارانشیم کبد تشخیص می‌دهد، نه از بافت کلی تصویر.
\end{quote}

اگر این فرض درست باشد، آن‌گاه بازنمایی درست مسئله نه حجم خام، بلکه بردار کوچکی از آماره‌های سطح‌اندام است. چنین بازنمایی‌ای سه مزیت دارد: ابعاد ورودی از چند میلیون وکسل به چند ده عدد کاهش می‌یابد و مسئله با داده‌ی آموزشی اندک قابل‌حل می‌شود؛ ویژگی‌ها واحد فیزیکی (\lr{HU}) دارند و در نتیجه میان دستگاه‌ها قابل‌انتقال‌ترند؛ و مدل حاصل \emph{تفسیرپذیر} است، زیرا سهم هر اندام در تصمیم قابل استخراج است.

\subsection*{استخراج ویژگی}

برای هر حجم، نقاب اندام با \lr{TotalSegmentator}~\cite{wasserthal2023} محاسبه و سپس برای هر اندام یک ویژگی—\emph{میانه‌ی} \lr{HU} درون آن اندام—استخراج می‌شود. انتخاب میانه به‌جای میانگین عمدی است: میانه در برابر داده‌ی پرت مقاوم است، و ضایعات کانونی، کلسیفیکاسیون و خطاهای موضعی بخش‌بندی دقیقاً همان داده‌ی پرتی هستند که میانگین را جابه‌جا می‌کنند.

نگاشت نام اندام به شناسه‌ی برچسب در زمان اجرا از نسخه‌ی نصب‌شده‌ی \lr{TotalSegmentator} خوانده می‌شود و در کد ثابت نشده است؛ در غیر این صورت، ارتقای نسخه‌ی این ابزار می‌توانست بی‌سروصدا ویژگی‌ها را جابه‌جا کند و مدل را بی‌آن‌که خطایی رخ دهد بی‌اعتبار سازد.

\subsection*{طبقه‌بند}

بردار ویژگی به یک طبقه‌بند تقویت گرادیان درختی\footnote{\lr{Gradient-boosted trees}} از نوع \lr{XGBoost}~\cite{chen2016xgboost} داده می‌شود. انتخاب این خانواده‌ی مدل، با توجه به ابعاد داده، طبیعی است: تعداد نمونه‌ها اندک و تعداد ویژگی‌ها کم است، ویژگی‌ها ناهمگن و دارای مقیاس فیزیکی‌اند، و رابطه‌ی میان روشنایی اندام و فاز آستانه‌ای و غیرخطی است—دقیقاً وضعیتی که درخت‌های تصمیم در آن بر شبکه‌های عصبی برتری دارند.

مدل با اعتبارسنجی متقاطع لایه‌ای\footnote{\lr{Stratified $k$-fold cross-validation}} ارزیابی می‌شود تا توزیع کلاس‌ها در هر تا حفظ شود. دو نسخه از مدل آموزش داده می‌شود: یکی روی مجموعه‌داده‌ی \lr{WAW-TACE} و دیگری روی \lr{VinDr-Multiphase}، تا انتقال‌پذیری میان مجموعه‌داده‌ها نیز سنجیده شود. نتایج در فصل~\ref{chap:experiments} گزارش می‌شود.

\section{ثبت تصویر بین‌فازی چندگذره}
\label{sec:m-reg}

\subsection*{ایده‌ی مرکزی: محدودسازی دامنه، نه وزن‌دهی}

معیارهای کلاسیک ثبت تصویر چندوجهی، به‌ویژه اطلاعات متقابل\footnote{\lr{Mutual information}}~\cite{wells1996,mattes2001}، بر این فرض بنا شده‌اند که میان شدت روشنایی دو تصویر یک رابطه‌ی آماری \emph{یکنواخت در سراسر تصویر} برقرار است. تزریق ماده‌ی حاجب این فرض را نقض می‌کند و نقض آن یکنواخت نیست: آئورت ممکن است بیش از صد واحد \lr{HU} روشن‌تر شود، کبد چند ده واحد، و استخوان اصلاً تغییر نکند. هیستوگرام مشترک میان دو فاز در نتیجه چندمُدی\footnote{\lr{Multi-modal}} می‌شود و معیاری که روی کل تصویر تعریف شده است، روی این ناهمگنی حاشیه‌گیری\footnote{\lr{Marginalise}} می‌کند و اطلاعات را از دست می‌دهد.

روش‌های آگاه از آناتومی موجود در ادبیات (بخش~\ref{sec:registration}) به این مشکل با \emph{وزن‌دهی} پاسخ می‌دهند: معیار همچنان روی کل تصویر محاسبه می‌شود، اما سهم نواحی مختلف در آن متفاوت است. پیشنهاد این پایان‌نامه از این نقطه جدا می‌شود:

\begin{quote}
\مهم{معیار تراز نباید در نواحی ناپایدار وزن کمتری بگیرد؛ نباید در آن نواحی \emph{محاسبه} شود.}
\end{quote}

به بیان صوری، اگر $\mathcal{M}$ نقاب اندام‌های پایدار در برابر فاز باشد، معیار $S$ نه به‌صورت $\sum_{x} w(x)\, s(x)$ با $w$ کوچک بیرون از $\mathcal{M}$، بلکه به‌صورت $\sum_{x \in \mathcal{M}} s(x)$ تعریف می‌شود. تفاوت این دو، در حالتی که $w$ به صفر میل کند، از نظر ریاضی از میان می‌رود؛ اما در عمل تفاوت اساسی است، زیرا در فرمول‌بندی وزنی، هیستوگرام مشترک—که اطلاعات متقابل از آن ساخته می‌شود—همچنان از وکسل‌های ناپایدار انباشته می‌شود و چندمُدی می‌ماند. محدودسازی دامنه، هیستوگرام را نیز محدود می‌کند.

\subsection*{ساخت نقاب چهارلایه}

نقاب $\mathcal{M}$ از حجم برچسب ادغام‌شده‌ی بخش~\ref{sec:m-data} ساخته می‌شود، با تقسیم اندام‌ها به لایه‌های زیر:

\begin{itemize}
\item \مهم{لایه‌ی ۱ — استخوان.} مهره‌ها، دنده‌ها، لگن، جناغ و طناب نخاعی. این ساختارها صلب‌اند و شدت روشنایی‌شان به‌کلی مستقل از فاز کنتراست است؛ بنابراین بالاترین اعتماد را دارند و مرجع صحت‌سنجی گذر صلب‌اند.
\item \مهم{لایه‌ی ۲ — بافت نرم پایدار.} کبد، طحال، کلیه‌ها، پانکراس، کیسه‌ی صفرا و عضلات محوری. این‌ها با کنتراست تغییر می‌کنند اما شکل و موقعیتشان پایدار است.
\item \مهم{لایه‌ی ۳ — سایر بافت نرم.} معده، مثانه، غدد فوق‌کلیه، لوب‌های ریه. تنها در نقاب \emph{ارزیابی} حضور دارند.
\item \مهم{حذف‌شده.} روده‌ی بزرگ و کوچک، دوازدهه و رکتوم. این ساختارها میان دو اسکن به‌دلیل حرکت دودی و گاز، جابه‌جایی غیرقابل‌مدل‌سازی دارند و حضورشان در معیار تراز، خطا وارد می‌کند نه اطلاعات.
\end{itemize}

نقاب ثبت از لایه‌های ۱ و ۲ ساخته می‌شود. پوشش نقاب (درصد وکسل‌های حجم) در هر مورد ثبت و پایش می‌شود؛ پوشش بسیار کم نشانه‌ی شکست بخش‌بندی است و باید پیش از تفسیر نتیجه‌ی ثبت تصویر شناسایی شود.

\subsection*{استانداردسازی پیش از ثبت}

پیش از هر محاسبه‌ی هندسی، هر سری از جهت‌های راهنمای \lr{DICOM} به قرارداد \lr{LPS} بازآرایی و به فاصله‌ی وکسل همسان‌گرد \lr{1.5} میلی‌متر بازنمونه‌برداری می‌شود. دو جزئیات این گام، که در نگاه نخست پیاده‌سازی به نظر می‌رسند، در عمل تعیین‌کننده‌اند.

نخست، \emph{مقدار پرکننده‌ی مرز} در بازنمونه‌برداری به‌صورت ثابت \lr{$-1024$} واحد \lr{HU} (هوا) تعیین می‌شود و نه کمینه‌ی خود تصویر. اسکنرهای سی‌تی وکسل‌های بیرون از میدان دید را با یک مقدار نشانه\footnote{\lr{Sentinel value}} ذخیره می‌کنند که پس از اعمال شیب و عرض‌ازمبدأ به \lr{$-4048$} واحد \lr{HU} می‌رسد؛ استفاده از این مقدار به‌عنوان پرکننده، هر وکسل مرزی حاصل از بازنمونه‌برداری را با عددی پر می‌کند که هیچ بافتی تولید نمی‌کند و هیستوگرام مشترک را که معیارهای مبتنی‌بر اطلاعات از آن ساخته می‌شوند آلوده می‌کند.

دوم، بازه‌ی \lr{HU} به \lr{$[-1024,\,+3000]$} برش داده می‌شود. کران پایین همان مقدار نشانه را حذف می‌کند. کران بالا عمداً بالاتر از \lr{$+1024$} انتخاب شده است: برش در \lr{$+1024$} استخوان متراکم قشری را می‌بُرد و کیفیت بخش‌بندی ساختارهای اسکلتی را---که در گذرهای بعدی نقش لنگر هندسی دارند---به‌طور محسوس کاهش می‌داد.

\subsection*{هم‌ترازسازی مرکز جرم با دروازه‌ی اطمینان}

دو سری از یک بررسی اغلب محدوده‌ی آناتومیک متفاوتی را پوشش می‌دهند. پیش از هر بهینه‌سازی مبتنی‌بر شدت روشنایی، جابه‌جایی درشت میان دو حجم از مرکز جرم سه‌بعدی چهار اندام \emph{لنگر}---کبد، طحال و دو کلیه---در فضای فیزیکی برآورد می‌شود. این چهار اندام بزرگ‌اند و موقعیتشان میان فازها پایدار است.

این گام سه‌گانه \emph{دروازه‌دار} است. یک نمره‌ی اطمینان $\gamma$، که سهم وزنی اندام‌های واقعاً یافت‌شده از کل وزن ممکن است، تعیین می‌کند که آیا برآورد قابل‌اتکاست؛ اگر $\gamma$ از \lr{0.40} کمتر باشد، یا اندازه‌ی جابه‌جایی از \lr{80} میلی‌متر بیشتر، یا پراکندگی برآوردهای تک‌اندام از \lr{12} میلی‌متر فراتر رود، هیچ جابه‌جایی‌ای اعمال نمی‌شود. منطق این محافظه‌کاری آن است که یک هم‌ترازسازی درشتِ نادرست بدتر از نبود آن است: بهینه‌سازی بعدی را به کمینه‌ی محلی اشتباه هدایت می‌کند.

افزون بر آن، دو ناحیه‌ی مرده\footnote{\lr{Deadband}} اعمال می‌شود: جابه‌جایی محور $z$ کمتر از \lr{7.5} میلی‌متر و جابه‌جایی درون‌صفحه‌ای کمتر از \lr{5} میلی‌متر نویز هم‌جلسه‌ای تلقی و نادیده گرفته می‌شوند، چرا که اعمال یک بازنمونه‌برداری اضافی برای تصحیحی در این مقیاس، بیش از آن‌که هم‌ترازی را بهبود دهد، وفاداری درون‌یابی را هزینه می‌کند.

\مهم{شاخه‌ی شاهد.} در کنار این مسیر، شاخه‌ی دومی اجرا می‌شود که \emph{همان} پنجره‌ی برش فیزیکی را می‌سازد اما هیچ هم‌ترازی‌ای اعمال نمی‌کند. چون تنها متغیر میان دو شاخه اعمالِ جابه‌جایی است، مقایسه‌ی آن‌ها سهم این گام را به‌صورت کنترل‌شده جدا می‌کند. شکل~\ref{fig:reg-pipeline} این انشعاب و جایگاه آن در کل زنجیره را نشان می‌دهد.

\subsection*{ثبت صلب سه‌گذره}

\paragraph{گذر صفر — مقداردهی اولیه‌ی بسته.}
به‌جای آغاز بهینه‌سازی از تبدیل همانی، یک تبدیل صلب اولیه به‌صورت بسته و بدون بهینه‌سازی تکراری محاسبه می‌شود. مرکز جرم هر اندام \emph{نقاب ثبت}---یعنی هم استخوان‌های لایه‌ی ۱ و هم بافت نرم پایدار لایه‌ی ۲---در دو حجم به‌عنوان یک زوج نقطه‌ی متناظر در نظر گرفته می‌شود و چرخش و انتقال بهینه میان دو مجموعه نقطه با تجزیه‌ی مقادیر منفرد\footnote{\lr{Singular value decomposition, SVD}} به روش کابش~\cite{kabsch1976} به دست می‌آید. یک نگهبان روی دترمینان ماتریس چرخش تضمین می‌کند که راه‌حل یک چرخش سره باشد و نه یک بازتاب؛ بدون این قید، دو حجم با جهت‌گیری آینه‌ای را می‌توان با نگاشت هر اندام به قرینه‌ی آناتومیک آن «هم‌تراز» کرد.

استفاده از هر دو لایه---و نه فقط استخوان---تصمیمی تجربی است: نسخه‌ی نخست این گذر تنها بر لنگرهای استخوانی تکیه می‌کرد و افزودن اندام‌های نرم پایدار به‌طور یکنواخت نتیجه‌ی بهتری داد، چرا که تعداد نقاط متناظر را افزایش و شرطی‌شدگی\footnote{\lr{Conditioning}} مسئله‌ی \lr{SVD} را بهبود می‌دهد. حاصل، یک نقطه‌ی شروع هندسی درست است که گذر بعدی را از دامنه‌ی جذب کمینه‌های محلی دور نگه می‌دارد.

\paragraph{گذر یک — پالایش با زیان آناتومیک.}
گذر دوم، تبدیل صلب را با کمینه‌سازی یک زیان \emph{آناتومیک} و نه یک زیان مبتنی‌بر شدت روشنایی پالایش می‌کند:
\begin{equation}
\label{eq:pass1-loss}
\mathcal{L}(T) \;=\; \lambda\, d_c(T) \;+\; \bigl(1 - \mathrm{Dice}_\epsilon(T)\bigr),
\qquad \lambda = 10 ,
\end{equation}
که در آن $d_c$ میانگین وزن‌دار جابه‌جایی مرکز جرم و $\mathrm{Dice}_\epsilon$ ضریب دایس \emph{فرسوده} است. هر دو جمله تنها روی یک زیرمجموعه‌ی گزیده از \عدد{۱۱} اندام محاسبه می‌شوند؛ ساختارهایی که موقعیتشان میان دو فاز ناپایدار است---معده و مری به‌دلیل حرکت دودی و پرشدگی، و عروق ایلیاک به‌دلیل تغییر شدید \lr{HU} با کنتراست---به‌صراحت از این دامنه بیرون گذاشته شده‌اند.

شعاع فرسایش به‌ازای هر اندام و متناسب با اندازه‌ی آن تعیین می‌شود (از \lr{6} میلی‌متر برای کبد تا \lr{1.5} میلی‌متر برای مهره‌ها)، تا ساختارهای کوچکِ دارای مرز نویزی نتوانند بر زیان مسلط شوند.

بهینه‌سازی با روش سیمپلکس نلدر-مید~\cite{nelder1965} انجام می‌گیرد. انتخاب یک بهینه‌ساز بدون مشتق در اینجا یک راه‌حل جایگزین نیست بلکه یک ضرورت است: ضریب دایس روی نقاب‌های دودویی گسسته محاسبه می‌شود و نسبت به پارامترهای تبدیل مشتق‌پذیر نیست، بنابراین هیچ بهینه‌ساز گرادیانی روی این زیان تعریف نمی‌شود.

\paragraph{گذر دو — پالایش ساختاری روی بزرگی گرادیان.}
گذر پایانی، معیار را از شدت روشنایی خام به بزرگی گرادیان\footnote{\lr{Gradient magnitude}} تغییر می‌دهد و آن را با همبستگی متقاطع نرمال‌شده‌ی محلی بیشینه می‌کند. منطق این تغییر آن است که تزریق ماده‌ی حاجب \emph{سطح} روشنایی را جابه‌جا می‌کند اما \emph{لبه}ها را در جای خود نگه می‌دارد؛ بنابراین معیاری که روی لبه تعریف شود، ذاتاً نسبت به فاز ناوردا\footnote{\lr{Invariant}}تر است. بزرگی گرادیان به \lr{200} واحد \lr{HU} بر میلی‌متر مقیدسازی\footnote{\lr{Clamp}} می‌شود تا مرزهای استخوان-هوا، که گرادیانی چند برابر بزرگ‌تر تولید می‌کنند، معیار را از مرزهای بافت نرم منحرف نکنند.

\paragraph{دروازه‌های پذیرش.}
گذرهای یک و دو هر یک با یک \emph{دروازه‌ی پذیرش} همراه‌اند: پس از اجرای گذر، خطای مرکز جرم دوباره اندازه‌گیری می‌شود و اگر نسبت به گذر پیشین بدتر شده باشد، تبدیل گذر پیشین حفظ و نتیجه‌ی گذر جدید دور ریخته می‌شود. این ساز و کار یک جزئیات مهندسی نیست بلکه یک تصمیم روش‌شناختی است: زنجیره‌ای از گذرهای پالایشی، بدون دروازه، هر شکست یک گذر را به همه‌ی گام‌های پایین‌دستی منتقل می‌کند و آن را در انبوه داده پنهان می‌سازد. با دروازه، بدترین حالتِ یک گذرِ ناموفق «بی‌اثر» است و نه «مخرب». چنان‌که در فصل~\ref{chap:experiments} نشان داده می‌شود، این تمایز در عمل تعیین‌کننده شد.

\subsection*{ثبت غیرصلب}

پالایش نهایی با یک تبدیل \lr{B}-اسپلاین آزادشکل~\cite{rueckert1999} روی شبکه‌ی کنترلی با فاصله‌ی \lr{25} میلی‌متر و بهینه‌ساز \lr{L-BFGS-B} انجام می‌شود. معیار، همانند گذرهای پیشین، به $\mathcal{M}$ محدود است. نقاب متحرک از نقاب ثابت به فضای تصویر متحرک بازنمونه‌برداری و با حاشیه‌ی اضافی گشاد\footnote{\lr{Dilate}} می‌شود تا خطای صلب باقی‌مانده را جذب کند. تبدیل روی \emph{تصاویر معیار} (نقشه‌ی لبه یا فاصله) برازش می‌شود اما بر تصویر سی‌تی اصلی اعمال می‌گردد.

میزان تاخوردگی\footnote{\lr{Folding}} میدان تغییرشکل، یعنی درصد وکسل‌هایی که دترمینان ژاکوبین در آن‌ها نامثبت است، در هر مورد \emph{گزارش} می‌شود اما به‌عنوان دروازه به کار نمی‌رود؛ گذاشتن دروازه بر آن، ثبت غیرصلب را بی‌سروصدا به ثبت صلب بازمی‌گرداند و همین کمیت است که بعداً امکان مقایسه‌ی منصفانه‌ی دقت و منظم‌بودن\footnote{\lr{Regularity}} میان روش‌ها را فراهم می‌کند.

پیاده‌سازی همه‌ی گذرها بر پایه‌ی \lr{SimpleITK}~\cite{lowekamp2013} است. شکل~\ref{fig:reg-pipeline} کل زنجیره را همراه با نقابی که در هر گام دامنه‌ی معیار را تعیین می‌کند خلاصه می‌کند.

% ---------------------------------------------------------------
%  شکل: نمودار جریان پایپ‌لاین پیش‌پردازش و ثبت تصویر بین‌فازی
%
%  نیازمند: tikz + کتابخانه‌های positioning, arrows.meta, calc
%           (در styles/common.tex و پیش از xepersian بارگذاری شده‌اند —
%            bidi باید آخرین بسته باشد، وگرنه خروجی شکل وارونه می‌شود)
%
%  استفاده:  % ---------------------------------------------------------------
%  شکل: نمودار جریان پایپ‌لاین پیش‌پردازش و ثبت تصویر بین‌فازی
%
%  نیازمند: tikz + کتابخانه‌های positioning, arrows.meta, calc
%           (در styles/common.tex و پیش از xepersian بارگذاری شده‌اند —
%            bidi باید آخرین بسته باشد، وگرنه خروجی شکل وارونه می‌شود)
%
%  استفاده:  % ---------------------------------------------------------------
%  شکل: نمودار جریان پایپ‌لاین پیش‌پردازش و ثبت تصویر بین‌فازی
%
%  نیازمند: tikz + کتابخانه‌های positioning, arrows.meta, calc
%           (در styles/common.tex و پیش از xepersian بارگذاری شده‌اند —
%            bidi باید آخرین بسته باشد، وگرنه خروجی شکل وارونه می‌شود)
%
%  استفاده:  \input{figs/pipeline_reg}
% ---------------------------------------------------------------

\begingroup
\newcommand{\Ftitle}{\scriptsize\bfseries}
\newcommand{\Fsub}{\scriptsize\color{xtire}}
\newcommand{\Fnum}{\scriptsize\color{xkhakestari}}

\begin{figure}[p]
\centering
\begin{tikzpicture}[
  font=\scriptsize,
  node distance=0mm,
  >={Stealth[length=1.7mm,width=1.4mm]},
  stage/.style={
    draw=xnokhodi, line width=.5pt, fill=white, rounded corners=1pt,
    text width=72mm, inner xsep=2.2mm, inner ysep=1.4mm, align=right},
  substage/.style={stage, text width=33mm},
  ctrl/.style={substage, fill=xkamrang, dash pattern=on 2pt off 1.4pt},
  maskbox/.style={
    draw=xabi, line width=.5pt, fill=xabikam, rounded corners=1pt,
    text width=33mm, inner xsep=2mm, inner ysep=1.3mm, align=right},
  dropbox/.style={maskbox, fill=white, draw=xnokhodi,
    dash pattern=on 2pt off 1.6pt},
  flow/.style={->, draw=xtire, line width=.6pt},
  plain/.style={draw=xtire, line width=.6pt},
  gov/.style={->, draw=xabi, line width=.45pt, dash pattern=on 1.2pt off 1.2pt},
  reject/.style={->, draw=xghermez, line width=.6pt, dash pattern=on 2.2pt off 1.5pt},
  accept/.style={->, draw=xsabz, line width=.6pt},
  edge/.style={font=\scriptsize, inner sep=1pt, fill=white},
  cnt/.style={font=\scriptsize, text=xkhakestari, inner sep=.6mm},
]

% ══════════════ ستون اصلی: مسیر داده ══════════════

\node[stage] (s1) {%
  {\Ftitle ۱ــ خواندن \lr{DICOM} و انتخاب سری}\\
  {\Fsub برای هر (بررسی، فاز) نازک‌ترین سری انتخاب می‌شود}\\
  {\Fnum \عدد{۲۶۵} بررسی \ \lr{$\cdot$}\ \عدد{۱۱۸۸} سری}};

\node[stage, below=3.2mm of s1] (s2) {%
  {\Ftitle ۲ــ استانداردسازی هندسی و شدت روشنایی}\\
  {\Fsub جهت \lr{LPS} \ \lr{$\cdot$}\ بازنمونه‌برداری همسان‌گرد \lr{1.5}\,میلی‌متر \ \lr{$\cdot$}\ پرکننده‌ی مرز \lr{$-1024$}\,\lr{HU}}\\
  {\Fsub برش بازه به \lr{$[-1024,+3000]$}\,\lr{HU}؛ حذف مقدار نشانه‌ی \lr{$-4048$} (\lr{$21.8\%$} وکسل‌ها)}\\
  {\Fnum \عدد{۷۷۲} حجم، همگی با فاصله‌ی وکسل یکسان}};

\node[stage, below=3.2mm of s2] (s3) {%
  {\Ftitle ۳ــ بخش‌بندی و ساخت نقاب چهارلایه}\\
  {\Fsub \lr{TotalSegmentator} و ادغام در یک حجم برچسب با نگاشت عددی ثابت}\\
  {\Fsub دروازه‌ی سلامت: دست‌کم \عدد{۲} از \عدد{۴} اندام لنگر و یک برچسب ستون فقرات}\\
  {\Fnum \عدد{۷۷۱} از \عدد{۷۷۲} حجم عبور می‌کند}};

% ---- انشعاب کنترل‌شده ----
\node[substage, below=6mm of s3.south west, anchor=north west,
      draw=xabi, line width=.7pt] (s4a) {%
  {\Ftitle ۴الف ــ برش + هم‌ترازی}\\
  {\Fsub جابه‌جایی مرکز جرم سه‌بعدی با}\\
  {\Fsub دروازه‌ی اطمینان \lr{$\gamma \ge 0.40$}}\\
  {\Fnum اعمال \lr{$z$}: \lr{$12.4\%$}، \lr{$xy$}: \lr{$8.8\%$}}};

\node[ctrl, right=6mm of s4a] (s4b) {%
  {\Ftitle ۴ب ــ فقط برش (شاهد)}\\
  {\Fsub همان پنجره‌ی برش فیزیکی،}\\
  {\Fsub بدون اعمال هیچ هم‌ترازی}\\
  {\Fnum سهم گام ۴الف را جدا می‌کند}};

\node[cnt, below=1mm of s4b.south east, anchor=north east]
     (n4) {\عدد{۱۶۹} بررسی دارای فاز بدون‌کنتراست \ \lr{$\cdot$}\ \عدد{۴۵۳} سری متحرک};

\node[stage, below=2.4mm of n4.south, anchor=north] (s5) {%
  {\Ftitle ۵ــ ثبت صلب، گذر صفر: کابش (\lr{SVD})}\\
  {\Fsub حل بسته‌ی چرخش و انتقال از مرکز جرم اندام‌های پایدار}\\
  {\Fnum بیشترین سهم در بافت پارانشیمی: دایس \lr{$0.742 \rightarrow 0.784$}}};

\node[stage, below=3.2mm of s5] (s6) {%
  {\Ftitle ۶ــ گذر یک: پالایش نلدر-مید}\\
  {\Fsub کمینه‌سازی \lr{$\lambda\, d_{c} + (1-\mathrm{Dice}_{\epsilon})$} با \lr{$\lambda=10$}، بدون نیاز به مشتق}\\
  {\Fnum بیشترین سهم در بافت اسکلتی: دایس \lr{$0.883 \rightarrow 0.908$}}};

\node[stage, below=9mm of s6] (s7) {%
  {\Ftitle ۷ــ گذر دو: همبستگی بزرگی گرادیان}\\
  {\Fsub \lr{$|\nabla \mathrm{HU}|$} با \lr{NCC} و هرم چندمقیاسی \lr{$[4,2,1]$}}\\
  {\Fnum پذیرفته‌شده در \عدد{۴} از \عدد{۶۹۹} ثبت}};

\node[stage, below=5.5mm of s7, fill=xkamrang] (s8) {%
  {\Ftitle ۸ــ ثبت غیرصلب \lr{B}-اسپلاین}\\
  {\Fsub شبکه‌ی کنترلی \lr{25}\,میلی‌متر \ \lr{$\cdot$}\ \lr{L-BFGS-B} \ \lr{$\cdot$}\ \عدد{۷} معیار قابل‌انتخاب}\\
  {\Fnum پیاده‌سازی شده؛ خروجی ارزیابی‌شده در دست نیست}};

\node[stage, below=3.2mm of s8] (s9) {%
  {\Ftitle ۹ــ ارزیابی روی نقاب کنارگذاشته}\\
  {\Fsub دایس فرسوده (\lr{6}\,میلی‌متر) \ \lr{$\cdot$}\ \lr{HD95} \ \lr{$\cdot$}\ جابه‌جایی مرکز جرم \ \lr{$\cdot$}\ \lr{$\%|J|\le 0$}}\\
  {\Fnum \عدد{۲۶} بررسی \ \lr{$\cdot$}\ \عدد{۵۵} جفت فاز \ \lr{$\cdot$}\ \عدد{۱۱} اندام}};

% ══════════════ ستون کناری: دامنه‌ی معیار ══════════════

\node[maskbox, left=6mm of s3, anchor=east] (m1) {%
  {\Ftitle نقاب کامل}\\
  {\Fsub لایه‌های ۱، ۲ و ۳. کادر برش و دامنه‌ی \emph{ارزیابی}.}};

\node[maskbox, left=6mm of s4a.west, anchor=east] (m2) {%
  {\Ftitle نقاب ثبت}\\
  {\Fsub لایه‌های ۱ و ۲: استخوان و بافت نرم پایدار. لنگر ۴الف، دامنه‌ی کابش، نقاب گذرهای ۷ و ۸.}};

\node[maskbox, left=6mm of s6, anchor=east] (m3) {%
  {\Ftitle زیرمجموعه‌ی پالایش}\\
  {\Fsub \عدد{۱۱} برچسب: کبد، طحال، کلیه‌ها، \lr{L1--L5}، طناب نخاعی، آئورت و بزرگ‌سیاهرگ زیرین. دامنه‌ی زیان گذر یک و همه‌ی دروازه‌ها.}};

\node[dropbox, left=6mm of s8, anchor=east] (m4) {%
  {\Ftitle لایه‌ی صفر ــ حذف کامل}\\
  {\Fsub روده، دوازدهه و رکتوم: حرکت دودی و گاز، سیگنال متناقض وارد معیار می‌کند.}};

% ══════════════ فلش‌های جریان ══════════════

\draw[flow] (s1) -- (s2);
\draw[flow] (s2) -- (s3);

\coordinate (fk) at ($(s3.south)+(0,-3mm)$);
\draw[plain] (s3.south) -- (fk);
\draw[plain] (s4a.north |- fk) -- (s4b.north |- fk);
\draw[flow]  (s4a.north |- fk) -- (s4a.north);
\draw[flow]  (s4b.north |- fk) -- (s4b.north);

\coordinate (jn) at ($(s5.north)+(0,2mm)$);
\draw[plain] (s4a.south) -- (s4a.south |- jn);
\draw[plain] (s4b.south) -- (s4b.south |- jn);
\draw[plain] (s4a.south |- jn) -- (s4b.south |- jn);
\draw[flow]  (s5.north |- jn) -- (s5.north);

\draw[flow] (s5) -- (s6);
\draw[flow] (s6) -- node[edge, text=xtire, pos=.66] {دروازه‌ی پذیرش} (s7);

% مسیر واقعی ۹۹٫۴ درصد ثبت‌ها: عبور از روی گذر دو
\coordinate (ba) at ($(s6.south)+(0,-1.4mm)$);
\coordinate (bb) at ($(s8.north)+(0,1.4mm)$);
\draw[reject]
  (ba) -- ($(ba)+(40mm,0)$) -- ($(bb)+(40mm,0)$) -- (bb) -- (s8.north);
\node[cnt, text=xghermez, align=right, anchor=west]
  at ($(ba)+(42mm,-13mm)$) {بازگشت به گذر یک\\ \lr{$695/699$}\\ \lr{$99.4\%$}};

\draw[accept] (s7.south) -- node[edge, text=xsabz] {\lr{$0.6\%$} پذیرش} (s8.north);
\draw[flow] (s8) -- (s9);

% ══════════════ فلش‌های حاکمیت نقاب ══════════════

% هر نقاب یک پیکان به مصرف‌کننده‌ی اصلی خود دارد؛ سایر گام‌هایی که همان
% نقاب را به کار می‌برند در متن جعبه فهرست شده‌اند تا نمودار شلوغ نشود.
\draw[gov] (m1.east) -- (s3.west);
\draw[gov] (m2.east) -- (s4a.west);
\draw[gov] (m3.east) -- (s6.west);
% دامنه‌ی ارزیابی: از نقاب کامل به گام ۹، از کانال آزاد سمت چپ
\coordinate (ch) at ($(m3.west)+(-4mm,0)$);
\draw[gov] (m1.west) -- (m1.west -| ch) -- (s9.west -| ch) -- (s9.west);

\end{tikzpicture}

\caption[نمودار جریان پایپ‌لاین ثبت تصویر بین‌فازی]{%
پایپ‌لاین پیش‌پردازش و ثبت تصویر بین‌فازی. پیکان‌های پیوسته مسیر داده‌اند و
پیکان‌های نقطه‌چین آبی نشان می‌دهند دامنه‌ی محاسبه‌ی معیار در هر گام کدام
نقاب است. شاخه‌ی ۴ب شاهدی کنترل‌شده است که همان پنجره‌ی برش را بدون
هم‌ترازی می‌سازد. مسیری که \عدد{۹۹٫۴} درصد ثبت‌ها در عمل طی می‌کنند پیکان
قرمز است و نه گذر دو.}
\label{fig:reg-pipeline}
\end{figure}
\endgroup

% ---------------------------------------------------------------

\begingroup
\newcommand{\Ftitle}{\scriptsize\bfseries}
\newcommand{\Fsub}{\scriptsize\color{xtire}}
\newcommand{\Fnum}{\scriptsize\color{xkhakestari}}

\begin{figure}[p]
\centering
\begin{tikzpicture}[
  font=\scriptsize,
  node distance=0mm,
  >={Stealth[length=1.7mm,width=1.4mm]},
  stage/.style={
    draw=xnokhodi, line width=.5pt, fill=white, rounded corners=1pt,
    text width=72mm, inner xsep=2.2mm, inner ysep=1.4mm, align=right},
  substage/.style={stage, text width=33mm},
  ctrl/.style={substage, fill=xkamrang, dash pattern=on 2pt off 1.4pt},
  maskbox/.style={
    draw=xabi, line width=.5pt, fill=xabikam, rounded corners=1pt,
    text width=33mm, inner xsep=2mm, inner ysep=1.3mm, align=right},
  dropbox/.style={maskbox, fill=white, draw=xnokhodi,
    dash pattern=on 2pt off 1.6pt},
  flow/.style={->, draw=xtire, line width=.6pt},
  plain/.style={draw=xtire, line width=.6pt},
  gov/.style={->, draw=xabi, line width=.45pt, dash pattern=on 1.2pt off 1.2pt},
  reject/.style={->, draw=xghermez, line width=.6pt, dash pattern=on 2.2pt off 1.5pt},
  accept/.style={->, draw=xsabz, line width=.6pt},
  edge/.style={font=\scriptsize, inner sep=1pt, fill=white},
  cnt/.style={font=\scriptsize, text=xkhakestari, inner sep=.6mm},
]

% ══════════════ ستون اصلی: مسیر داده ══════════════

\node[stage] (s1) {%
  {\Ftitle ۱ــ خواندن \lr{DICOM} و انتخاب سری}\\
  {\Fsub برای هر (بررسی، فاز) نازک‌ترین سری انتخاب می‌شود}\\
  {\Fnum \عدد{۲۶۵} بررسی \ \lr{$\cdot$}\ \عدد{۱۱۸۸} سری}};

\node[stage, below=3.2mm of s1] (s2) {%
  {\Ftitle ۲ــ استانداردسازی هندسی و شدت روشنایی}\\
  {\Fsub جهت \lr{LPS} \ \lr{$\cdot$}\ بازنمونه‌برداری همسان‌گرد \lr{1.5}\,میلی‌متر \ \lr{$\cdot$}\ پرکننده‌ی مرز \lr{$-1024$}\,\lr{HU}}\\
  {\Fsub برش بازه به \lr{$[-1024,+3000]$}\,\lr{HU}؛ حذف مقدار نشانه‌ی \lr{$-4048$} (\lr{$21.8\%$} وکسل‌ها)}\\
  {\Fnum \عدد{۷۷۲} حجم، همگی با فاصله‌ی وکسل یکسان}};

\node[stage, below=3.2mm of s2] (s3) {%
  {\Ftitle ۳ــ بخش‌بندی و ساخت نقاب چهارلایه}\\
  {\Fsub \lr{TotalSegmentator} و ادغام در یک حجم برچسب با نگاشت عددی ثابت}\\
  {\Fsub دروازه‌ی سلامت: دست‌کم \عدد{۲} از \عدد{۴} اندام لنگر و یک برچسب ستون فقرات}\\
  {\Fnum \عدد{۷۷۱} از \عدد{۷۷۲} حجم عبور می‌کند}};

% ---- انشعاب کنترل‌شده ----
\node[substage, below=6mm of s3.south west, anchor=north west,
      draw=xabi, line width=.7pt] (s4a) {%
  {\Ftitle ۴الف ــ برش + هم‌ترازی}\\
  {\Fsub جابه‌جایی مرکز جرم سه‌بعدی با}\\
  {\Fsub دروازه‌ی اطمینان \lr{$\gamma \ge 0.40$}}\\
  {\Fnum اعمال \lr{$z$}: \lr{$12.4\%$}، \lr{$xy$}: \lr{$8.8\%$}}};

\node[ctrl, right=6mm of s4a] (s4b) {%
  {\Ftitle ۴ب ــ فقط برش (شاهد)}\\
  {\Fsub همان پنجره‌ی برش فیزیکی،}\\
  {\Fsub بدون اعمال هیچ هم‌ترازی}\\
  {\Fnum سهم گام ۴الف را جدا می‌کند}};

\node[cnt, below=1mm of s4b.south east, anchor=north east]
     (n4) {\عدد{۱۶۹} بررسی دارای فاز بدون‌کنتراست \ \lr{$\cdot$}\ \عدد{۴۵۳} سری متحرک};

\node[stage, below=2.4mm of n4.south, anchor=north] (s5) {%
  {\Ftitle ۵ــ ثبت صلب، گذر صفر: کابش (\lr{SVD})}\\
  {\Fsub حل بسته‌ی چرخش و انتقال از مرکز جرم اندام‌های پایدار}\\
  {\Fnum بیشترین سهم در بافت پارانشیمی: دایس \lr{$0.742 \rightarrow 0.784$}}};

\node[stage, below=3.2mm of s5] (s6) {%
  {\Ftitle ۶ــ گذر یک: پالایش نلدر-مید}\\
  {\Fsub کمینه‌سازی \lr{$\lambda\, d_{c} + (1-\mathrm{Dice}_{\epsilon})$} با \lr{$\lambda=10$}، بدون نیاز به مشتق}\\
  {\Fnum بیشترین سهم در بافت اسکلتی: دایس \lr{$0.883 \rightarrow 0.908$}}};

\node[stage, below=9mm of s6] (s7) {%
  {\Ftitle ۷ــ گذر دو: همبستگی بزرگی گرادیان}\\
  {\Fsub \lr{$|\nabla \mathrm{HU}|$} با \lr{NCC} و هرم چندمقیاسی \lr{$[4,2,1]$}}\\
  {\Fnum پذیرفته‌شده در \عدد{۴} از \عدد{۶۹۹} ثبت}};

\node[stage, below=5.5mm of s7, fill=xkamrang] (s8) {%
  {\Ftitle ۸ــ ثبت غیرصلب \lr{B}-اسپلاین}\\
  {\Fsub شبکه‌ی کنترلی \lr{25}\,میلی‌متر \ \lr{$\cdot$}\ \lr{L-BFGS-B} \ \lr{$\cdot$}\ \عدد{۷} معیار قابل‌انتخاب}\\
  {\Fnum پیاده‌سازی شده؛ خروجی ارزیابی‌شده در دست نیست}};

\node[stage, below=3.2mm of s8] (s9) {%
  {\Ftitle ۹ــ ارزیابی روی نقاب کنارگذاشته}\\
  {\Fsub دایس فرسوده (\lr{6}\,میلی‌متر) \ \lr{$\cdot$}\ \lr{HD95} \ \lr{$\cdot$}\ جابه‌جایی مرکز جرم \ \lr{$\cdot$}\ \lr{$\%|J|\le 0$}}\\
  {\Fnum \عدد{۲۶} بررسی \ \lr{$\cdot$}\ \عدد{۵۵} جفت فاز \ \lr{$\cdot$}\ \عدد{۱۱} اندام}};

% ══════════════ ستون کناری: دامنه‌ی معیار ══════════════

\node[maskbox, left=6mm of s3, anchor=east] (m1) {%
  {\Ftitle نقاب کامل}\\
  {\Fsub لایه‌های ۱، ۲ و ۳. کادر برش و دامنه‌ی \emph{ارزیابی}.}};

\node[maskbox, left=6mm of s4a.west, anchor=east] (m2) {%
  {\Ftitle نقاب ثبت}\\
  {\Fsub لایه‌های ۱ و ۲: استخوان و بافت نرم پایدار. لنگر ۴الف، دامنه‌ی کابش، نقاب گذرهای ۷ و ۸.}};

\node[maskbox, left=6mm of s6, anchor=east] (m3) {%
  {\Ftitle زیرمجموعه‌ی پالایش}\\
  {\Fsub \عدد{۱۱} برچسب: کبد، طحال، کلیه‌ها، \lr{L1--L5}، طناب نخاعی، آئورت و بزرگ‌سیاهرگ زیرین. دامنه‌ی زیان گذر یک و همه‌ی دروازه‌ها.}};

\node[dropbox, left=6mm of s8, anchor=east] (m4) {%
  {\Ftitle لایه‌ی صفر ــ حذف کامل}\\
  {\Fsub روده، دوازدهه و رکتوم: حرکت دودی و گاز، سیگنال متناقض وارد معیار می‌کند.}};

% ══════════════ فلش‌های جریان ══════════════

\draw[flow] (s1) -- (s2);
\draw[flow] (s2) -- (s3);

\coordinate (fk) at ($(s3.south)+(0,-3mm)$);
\draw[plain] (s3.south) -- (fk);
\draw[plain] (s4a.north |- fk) -- (s4b.north |- fk);
\draw[flow]  (s4a.north |- fk) -- (s4a.north);
\draw[flow]  (s4b.north |- fk) -- (s4b.north);

\coordinate (jn) at ($(s5.north)+(0,2mm)$);
\draw[plain] (s4a.south) -- (s4a.south |- jn);
\draw[plain] (s4b.south) -- (s4b.south |- jn);
\draw[plain] (s4a.south |- jn) -- (s4b.south |- jn);
\draw[flow]  (s5.north |- jn) -- (s5.north);

\draw[flow] (s5) -- (s6);
\draw[flow] (s6) -- node[edge, text=xtire, pos=.66] {دروازه‌ی پذیرش} (s7);

% مسیر واقعی ۹۹٫۴ درصد ثبت‌ها: عبور از روی گذر دو
\coordinate (ba) at ($(s6.south)+(0,-1.4mm)$);
\coordinate (bb) at ($(s8.north)+(0,1.4mm)$);
\draw[reject]
  (ba) -- ($(ba)+(40mm,0)$) -- ($(bb)+(40mm,0)$) -- (bb) -- (s8.north);
\node[cnt, text=xghermez, align=right, anchor=west]
  at ($(ba)+(42mm,-13mm)$) {بازگشت به گذر یک\\ \lr{$695/699$}\\ \lr{$99.4\%$}};

\draw[accept] (s7.south) -- node[edge, text=xsabz] {\lr{$0.6\%$} پذیرش} (s8.north);
\draw[flow] (s8) -- (s9);

% ══════════════ فلش‌های حاکمیت نقاب ══════════════

% هر نقاب یک پیکان به مصرف‌کننده‌ی اصلی خود دارد؛ سایر گام‌هایی که همان
% نقاب را به کار می‌برند در متن جعبه فهرست شده‌اند تا نمودار شلوغ نشود.
\draw[gov] (m1.east) -- (s3.west);
\draw[gov] (m2.east) -- (s4a.west);
\draw[gov] (m3.east) -- (s6.west);
% دامنه‌ی ارزیابی: از نقاب کامل به گام ۹، از کانال آزاد سمت چپ
\coordinate (ch) at ($(m3.west)+(-4mm,0)$);
\draw[gov] (m1.west) -- (m1.west -| ch) -- (s9.west -| ch) -- (s9.west);

\end{tikzpicture}

\caption[نمودار جریان پایپ‌لاین ثبت تصویر بین‌فازی]{%
پایپ‌لاین پیش‌پردازش و ثبت تصویر بین‌فازی. پیکان‌های پیوسته مسیر داده‌اند و
پیکان‌های نقطه‌چین آبی نشان می‌دهند دامنه‌ی محاسبه‌ی معیار در هر گام کدام
نقاب است. شاخه‌ی ۴ب شاهدی کنترل‌شده است که همان پنجره‌ی برش را بدون
هم‌ترازی می‌سازد. مسیری که \عدد{۹۹٫۴} درصد ثبت‌ها در عمل طی می‌کنند پیکان
قرمز است و نه گذر دو.}
\label{fig:reg-pipeline}
\end{figure}
\endgroup

% ---------------------------------------------------------------

\begingroup
\newcommand{\Ftitle}{\scriptsize\bfseries}
\newcommand{\Fsub}{\scriptsize\color{xtire}}
\newcommand{\Fnum}{\scriptsize\color{xkhakestari}}

\begin{figure}[p]
\centering
\begin{tikzpicture}[
  font=\scriptsize,
  node distance=0mm,
  >={Stealth[length=1.7mm,width=1.4mm]},
  stage/.style={
    draw=xnokhodi, line width=.5pt, fill=white, rounded corners=1pt,
    text width=72mm, inner xsep=2.2mm, inner ysep=1.4mm, align=right},
  substage/.style={stage, text width=33mm},
  ctrl/.style={substage, fill=xkamrang, dash pattern=on 2pt off 1.4pt},
  maskbox/.style={
    draw=xabi, line width=.5pt, fill=xabikam, rounded corners=1pt,
    text width=33mm, inner xsep=2mm, inner ysep=1.3mm, align=right},
  dropbox/.style={maskbox, fill=white, draw=xnokhodi,
    dash pattern=on 2pt off 1.6pt},
  flow/.style={->, draw=xtire, line width=.6pt},
  plain/.style={draw=xtire, line width=.6pt},
  gov/.style={->, draw=xabi, line width=.45pt, dash pattern=on 1.2pt off 1.2pt},
  reject/.style={->, draw=xghermez, line width=.6pt, dash pattern=on 2.2pt off 1.5pt},
  accept/.style={->, draw=xsabz, line width=.6pt},
  edge/.style={font=\scriptsize, inner sep=1pt, fill=white},
  cnt/.style={font=\scriptsize, text=xkhakestari, inner sep=.6mm},
]

% ══════════════ ستون اصلی: مسیر داده ══════════════

\node[stage] (s1) {%
  {\Ftitle ۱ــ خواندن \lr{DICOM} و انتخاب سری}\\
  {\Fsub برای هر (بررسی، فاز) نازک‌ترین سری انتخاب می‌شود}\\
  {\Fnum \عدد{۲۶۵} بررسی \ \lr{$\cdot$}\ \عدد{۱۱۸۸} سری}};

\node[stage, below=3.2mm of s1] (s2) {%
  {\Ftitle ۲ــ استانداردسازی هندسی و شدت روشنایی}\\
  {\Fsub جهت \lr{LPS} \ \lr{$\cdot$}\ بازنمونه‌برداری همسان‌گرد \lr{1.5}\,میلی‌متر \ \lr{$\cdot$}\ پرکننده‌ی مرز \lr{$-1024$}\,\lr{HU}}\\
  {\Fsub برش بازه به \lr{$[-1024,+3000]$}\,\lr{HU}؛ حذف مقدار نشانه‌ی \lr{$-4048$} (\lr{$21.8\%$} وکسل‌ها)}\\
  {\Fnum \عدد{۷۷۲} حجم، همگی با فاصله‌ی وکسل یکسان}};

\node[stage, below=3.2mm of s2] (s3) {%
  {\Ftitle ۳ــ بخش‌بندی و ساخت نقاب چهارلایه}\\
  {\Fsub \lr{TotalSegmentator} و ادغام در یک حجم برچسب با نگاشت عددی ثابت}\\
  {\Fsub دروازه‌ی سلامت: دست‌کم \عدد{۲} از \عدد{۴} اندام لنگر و یک برچسب ستون فقرات}\\
  {\Fnum \عدد{۷۷۱} از \عدد{۷۷۲} حجم عبور می‌کند}};

% ---- انشعاب کنترل‌شده ----
\node[substage, below=6mm of s3.south west, anchor=north west,
      draw=xabi, line width=.7pt] (s4a) {%
  {\Ftitle ۴الف ــ برش + هم‌ترازی}\\
  {\Fsub جابه‌جایی مرکز جرم سه‌بعدی با}\\
  {\Fsub دروازه‌ی اطمینان \lr{$\gamma \ge 0.40$}}\\
  {\Fnum اعمال \lr{$z$}: \lr{$12.4\%$}، \lr{$xy$}: \lr{$8.8\%$}}};

\node[ctrl, right=6mm of s4a] (s4b) {%
  {\Ftitle ۴ب ــ فقط برش (شاهد)}\\
  {\Fsub همان پنجره‌ی برش فیزیکی،}\\
  {\Fsub بدون اعمال هیچ هم‌ترازی}\\
  {\Fnum سهم گام ۴الف را جدا می‌کند}};

\node[cnt, below=1mm of s4b.south east, anchor=north east]
     (n4) {\عدد{۱۶۹} بررسی دارای فاز بدون‌کنتراست \ \lr{$\cdot$}\ \عدد{۴۵۳} سری متحرک};

\node[stage, below=2.4mm of n4.south, anchor=north] (s5) {%
  {\Ftitle ۵ــ ثبت صلب، گذر صفر: کابش (\lr{SVD})}\\
  {\Fsub حل بسته‌ی چرخش و انتقال از مرکز جرم اندام‌های پایدار}\\
  {\Fnum بیشترین سهم در بافت پارانشیمی: دایس \lr{$0.742 \rightarrow 0.784$}}};

\node[stage, below=3.2mm of s5] (s6) {%
  {\Ftitle ۶ــ گذر یک: پالایش نلدر-مید}\\
  {\Fsub کمینه‌سازی \lr{$\lambda\, d_{c} + (1-\mathrm{Dice}_{\epsilon})$} با \lr{$\lambda=10$}، بدون نیاز به مشتق}\\
  {\Fnum بیشترین سهم در بافت اسکلتی: دایس \lr{$0.883 \rightarrow 0.908$}}};

\node[stage, below=9mm of s6] (s7) {%
  {\Ftitle ۷ــ گذر دو: همبستگی بزرگی گرادیان}\\
  {\Fsub \lr{$|\nabla \mathrm{HU}|$} با \lr{NCC} و هرم چندمقیاسی \lr{$[4,2,1]$}}\\
  {\Fnum پذیرفته‌شده در \عدد{۴} از \عدد{۶۹۹} ثبت}};

\node[stage, below=5.5mm of s7, fill=xkamrang] (s8) {%
  {\Ftitle ۸ــ ثبت غیرصلب \lr{B}-اسپلاین}\\
  {\Fsub شبکه‌ی کنترلی \lr{25}\,میلی‌متر \ \lr{$\cdot$}\ \lr{L-BFGS-B} \ \lr{$\cdot$}\ \عدد{۷} معیار قابل‌انتخاب}\\
  {\Fnum پیاده‌سازی شده؛ خروجی ارزیابی‌شده در دست نیست}};

\node[stage, below=3.2mm of s8] (s9) {%
  {\Ftitle ۹ــ ارزیابی روی نقاب کنارگذاشته}\\
  {\Fsub دایس فرسوده (\lr{6}\,میلی‌متر) \ \lr{$\cdot$}\ \lr{HD95} \ \lr{$\cdot$}\ جابه‌جایی مرکز جرم \ \lr{$\cdot$}\ \lr{$\%|J|\le 0$}}\\
  {\Fnum \عدد{۲۶} بررسی \ \lr{$\cdot$}\ \عدد{۵۵} جفت فاز \ \lr{$\cdot$}\ \عدد{۱۱} اندام}};

% ══════════════ ستون کناری: دامنه‌ی معیار ══════════════

\node[maskbox, left=6mm of s3, anchor=east] (m1) {%
  {\Ftitle نقاب کامل}\\
  {\Fsub لایه‌های ۱، ۲ و ۳. کادر برش و دامنه‌ی \emph{ارزیابی}.}};

\node[maskbox, left=6mm of s4a.west, anchor=east] (m2) {%
  {\Ftitle نقاب ثبت}\\
  {\Fsub لایه‌های ۱ و ۲: استخوان و بافت نرم پایدار. لنگر ۴الف، دامنه‌ی کابش، نقاب گذرهای ۷ و ۸.}};

\node[maskbox, left=6mm of s6, anchor=east] (m3) {%
  {\Ftitle زیرمجموعه‌ی پالایش}\\
  {\Fsub \عدد{۱۱} برچسب: کبد، طحال، کلیه‌ها، \lr{L1--L5}، طناب نخاعی، آئورت و بزرگ‌سیاهرگ زیرین. دامنه‌ی زیان گذر یک و همه‌ی دروازه‌ها.}};

\node[dropbox, left=6mm of s8, anchor=east] (m4) {%
  {\Ftitle لایه‌ی صفر ــ حذف کامل}\\
  {\Fsub روده، دوازدهه و رکتوم: حرکت دودی و گاز، سیگنال متناقض وارد معیار می‌کند.}};

% ══════════════ فلش‌های جریان ══════════════

\draw[flow] (s1) -- (s2);
\draw[flow] (s2) -- (s3);

\coordinate (fk) at ($(s3.south)+(0,-3mm)$);
\draw[plain] (s3.south) -- (fk);
\draw[plain] (s4a.north |- fk) -- (s4b.north |- fk);
\draw[flow]  (s4a.north |- fk) -- (s4a.north);
\draw[flow]  (s4b.north |- fk) -- (s4b.north);

\coordinate (jn) at ($(s5.north)+(0,2mm)$);
\draw[plain] (s4a.south) -- (s4a.south |- jn);
\draw[plain] (s4b.south) -- (s4b.south |- jn);
\draw[plain] (s4a.south |- jn) -- (s4b.south |- jn);
\draw[flow]  (s5.north |- jn) -- (s5.north);

\draw[flow] (s5) -- (s6);
\draw[flow] (s6) -- node[edge, text=xtire, pos=.66] {دروازه‌ی پذیرش} (s7);

% مسیر واقعی ۹۹٫۴ درصد ثبت‌ها: عبور از روی گذر دو
\coordinate (ba) at ($(s6.south)+(0,-1.4mm)$);
\coordinate (bb) at ($(s8.north)+(0,1.4mm)$);
\draw[reject]
  (ba) -- ($(ba)+(40mm,0)$) -- ($(bb)+(40mm,0)$) -- (bb) -- (s8.north);
\node[cnt, text=xghermez, align=right, anchor=west]
  at ($(ba)+(42mm,-13mm)$) {بازگشت به گذر یک\\ \lr{$695/699$}\\ \lr{$99.4\%$}};

\draw[accept] (s7.south) -- node[edge, text=xsabz] {\lr{$0.6\%$} پذیرش} (s8.north);
\draw[flow] (s8) -- (s9);

% ══════════════ فلش‌های حاکمیت نقاب ══════════════

% هر نقاب یک پیکان به مصرف‌کننده‌ی اصلی خود دارد؛ سایر گام‌هایی که همان
% نقاب را به کار می‌برند در متن جعبه فهرست شده‌اند تا نمودار شلوغ نشود.
\draw[gov] (m1.east) -- (s3.west);
\draw[gov] (m2.east) -- (s4a.west);
\draw[gov] (m3.east) -- (s6.west);
% دامنه‌ی ارزیابی: از نقاب کامل به گام ۹، از کانال آزاد سمت چپ
\coordinate (ch) at ($(m3.west)+(-4mm,0)$);
\draw[gov] (m1.west) -- (m1.west -| ch) -- (s9.west -| ch) -- (s9.west);

\end{tikzpicture}

\caption[نمودار جریان پایپ‌لاین ثبت تصویر بین‌فازی]{%
پایپ‌لاین پیش‌پردازش و ثبت تصویر بین‌فازی. پیکان‌های پیوسته مسیر داده‌اند و
پیکان‌های نقطه‌چین آبی نشان می‌دهند دامنه‌ی محاسبه‌ی معیار در هر گام کدام
نقاب است. شاخه‌ی ۴ب شاهدی کنترل‌شده است که همان پنجره‌ی برش را بدون
هم‌ترازی می‌سازد. مسیری که \عدد{۹۹٫۴} درصد ثبت‌ها در عمل طی می‌کنند پیکان
قرمز است و نه گذر دو.}
\label{fig:reg-pipeline}
\end{figure}
\endgroup


\subsection*{مجموعه‌ی معیارهای قابل‌انتخاب}

برای آن‌که بتوان سهم \emph{محدودسازی دامنه} را از سهم \emph{انتخاب معیار} جدا کرد، گذر غیرصلب با هفت معیار متفاوت قابل اجراست: اطلاعات متقابل ماتس روی \lr{HU} خام، همبستگی متقاطع نرمال‌شده روی بزرگی گرادیان مرکزی، همان معیار روی گرادیان هموارشده‌ی سوبل، نقشه‌ی لبه‌ی دودویی‌شده با آستانه‌ی اتسو، فاصله‌ی علامت‌دار نقاب اندام، توصیفگر ساختاری \lr{MIND}~\cite{heinrich2012}، و ترکیب ضربی \lr{MIND} با لبه‌ی دودویی. علاوه بر این، هر یک از این معیارها را می‌توان با دامنه‌ی محدودشده به $\mathcal{M}$ یا روی کل تصویر اجرا کرد. حاصل، یک طرح عاملی\footnote{\lr{Factorial design}} دوبعدی (معیار $\times$ دامنه) است که آزمون مستقیم ادعای مرکزی این بخش را ممکن می‌سازد؛ نتایج آن در فصل~\ref{chap:experiments} می‌آید.

\subsection*{معیارهای ارزیابی ثبت تصویر}

کیفیت هم‌ترازی با چهار کمیت سنجیده می‌شود که همگی روی نقاب \emph{ارزیابی} و نه نقاب ثبت محاسبه می‌شوند—در غیر این صورت معیار ارزیابی همان چیزی می‌بود که بهینه شده است:

\begin{itemize}
\item \مهم{ضریب دایس فرسوده\footnote{\lr{Eroded Dice}}}: ضریب دایس پس از فرسایش هر دو نقاب به‌اندازه‌ی \lr{6} میلی‌متر در فضای فیزیکی. فرسایش، حساسیت معیار به نویز مرزی بخش‌بندی را کاهش می‌دهد و آن را به تراز \emph{مرکز} اندام حساس‌تر می‌کند.
\item \مهم{فاصله‌ی هاسدورف صدک ۹۵}: کران بالای فاصله‌ی سطحی، مقاوم نسبت به تعداد اندکی وکسل پرت.
\item \مهم{جابه‌جایی مرکز جرم}: بر حسب میلی‌متر، روی نقاب‌های فرسوده‌نشده.
\item \مهم{همبستگی گرادیان درون اندام}: میزان انطباق ساختار لبه‌ای، مستقل از سطح روشنایی.
\item \مهم{همبستگی توصیفگر \lr{MIND}}: انطباق ساختار خودمشابه، معیاری مستقل از هر دو مورد بالا.
\item \مهم{درصد تاخوردگی}: تنها برای روش‌های غیرصلب؛ سهم وکسل‌هایی با دترمینان ژاکوبین نامثبت. این کمیت \emph{دقت} را نمی‌سنجد بلکه \emph{اعتبار فیزیکی} تغییرشکل را می‌سنجد و بدون آن، مقایسه‌ی یک روش غیرصلب با یک روش صلب بی‌معناست.
\end{itemize}

همبستگی متقاطع روی \lr{HU} خام عمداً از این فهرست کنار گذاشته شده است: این معیار وجود رابطه‌ی خطی میان شدت روشنایی متناظر دو تصویر را فرض می‌کند و تزریق ماده‌ی حاجب دقیقاً همین فرض را نقض می‌کند، بنابراین مقدار آن میان دو فاز نویز است و نه سیگنال.

اندام‌ها برای گزارش به سه گروه پارانشیمی، عروقی و اسکلتی دسته‌بندی می‌شوند. گروه اسکلتی نقش \emph{شاهد منفی} دارد: چون استخوان با کنتراست تغییر نمی‌کند، خطای باقی‌مانده در آن نشانه‌ی شکست ثبت تصویر است و نه دشواری آناتومیک.

\section{سنتز فاز کنتراست‌دار از فاز بدون‌کنتراست}
\label{sec:m-syn}

\subsection*{صورت‌بندی مسئله}

جفت‌های هم‌ترازشده‌ی حاصل از بخش~\ref{sec:m-reg} یک مسئله‌ی ترجمه‌ی تصویر-به-تصویر نظارت‌شده تعریف می‌کنند: با داشتن حجم بدون‌کنتراست $x$، حجم کنتراست‌دار $y$ همان بیمار در همان هندسه بازسازی شود. صورت‌بندی متعارف، کمینه‌سازی یک زیان بازسازی $\mathcal{L}(G(x), y)$ است.

این پایان‌نامه سه خانواده‌ی مدل را روی همین صورت‌بندی و همین تقسیم داده مقایسه می‌کند: مدل قطعی مبتنی‌بر شبکه‌ی مولد-متخاصم، مدل انتشار\footnote{\lr{Diffusion model}}، و رگرسیون ناهم‌واریانس\footnote{\lr{Heteroscedastic regression}}. منطق این انتخاب سه‌گانه در بخش~\ref{sec:m-uncertainty} روشن می‌شود.

\subsection*{مولد قطعی}

معماری پایه یک شبکه‌ی \lr{U-Net}~\cite{ronneberger2015unet} دوبعدی است که روی وصله‌های\footnote{\lr{Patch}} $128 \times 128$ از برش‌های محوری آموزش می‌بیند. آموزش وصله‌ای—به‌جای برش کامل—حافظه را آزاد و تعداد نمونه‌های مؤثر را زیاد می‌کند.

\paragraph{نرمال‌سازی.}
لایه‌های نرمال‌سازی گروهی~\cite{wu2018group} به‌جای نرمال‌سازی نمونه‌ای~\cite{ulyanov2016instance} به کار می‌روند. دلیل این انتخاب در فصل~\ref{chap:experiments} به‌صورت کمّی ارائه می‌شود و به مسئله‌ی درز\footnote{\lr{Seam}} میان وصله‌ها در استنتاج کاشی‌کاری‌شده مربوط است.

\paragraph{متمایزکننده.}
یک متمایزکننده‌ی وصله‌ای\footnote{\lr{PatchGAN}}~\cite{isola2017image} به‌صورت \emph{شرطی} به کار می‌رود: ورودی آن الحاق تصویر مبدأ و تصویر مقصد است، نه تصویر مقصد به‌تنهایی. متمایزکننده‌ی غیرشرطی نمی‌تواند تشخیص دهد که تصویر تولیدشده به \emph{این} ورودی مربوط است یا نه، و صرفاً واقع‌نمایی کلی را پاداش می‌دهد.

\paragraph{تابع زیان.}
زیان کل ترکیبی است از زیان $L_1$ سراسری، یک زیان $L_1$ وزن‌دار در سطح اندام، و جمله‌ی متخاصم. وزن‌های سطح‌اندام لایه‌بندی‌شده‌اند و بر اساس اهمیت هر اندام در تعیین فاز کنتراست تعیین می‌شوند. روده وزن \emph{صفر} می‌گیرد: محتوای گاز و مدفوع آن میان دو اسکن تغییر می‌کند و از روی فاز بدون‌کنتراست قابل استنتاج نیست، پس آموزش مدل برای بازسازی آن، ظرفیت را صرف نویز می‌کند. روده همچنان \emph{ارزیابی} می‌شود، صرفاً \emph{آموزش} داده نمی‌شود—تمایزی که در فصل~\ref{chap:experiments} به‌صراحت دفاع می‌شود.

\paragraph{بهینه‌سازی.}
نرخ یادگیری از زمان‌بندی کسینوسی با راه‌اندازی مجدد گرم~\cite{loshchilov2017sgdr} پیروی می‌کند. وزن جمله‌ی $L_1$ در طول آموزش کاهش می‌یابد تا در مراحل اولیه ساختار کلی و در مراحل بعدی جزئیات بافتی شکل بگیرد.

\subsection*{زمینه‌ی محور $z$: مدل ۲٫۵ بعدی}

مدل پایه هر برش را مستقل پردازش می‌کند و از این رو به پیوستگی میان‌برشی بی‌اعتناست. برای افزودن زمینه‌ی محور $z$ بدون هزینه‌ی یک شبکه‌ی سه‌بعدی کامل، یک پیکربندی ۲٫۵ بعدی به کار می‌رود: $2k+1$ برش مجاور به‌عنوان کانال‌های ورودی داده می‌شوند و تنها برش میانی پیش‌بینی می‌شود. شبکه دوبعدی باقی می‌ماند و تعداد پارامترها تقریباً تغییر نمی‌کند. پنجره‌ی $z$ در مرزهای حجم بریده\footnote{\lr{Edge-clamped}} می‌شود تا رفتار آموزش و استنتاج در ابتدا و انتهای حجم یکسان بماند و هیچ برشی از دست نرود.

\subsection*{شرطی‌سازی}

دو نوع شرطی‌سازی با سازوکار مدولاسیون ویژگی خطی\footnote{\lr{Feature-wise Linear Modulation, FiLM}}~\cite{perez2018film} پیاده شده است:

\begin{itemize}
\item \مهم{شرطی‌سازی فاز}: یک بردار که فاز هدف (شریانی یا وریدی) را مشخص می‌کند و به یک مدل واحد اجازه می‌دهد هر دو فاز را تولید کند.
\item \مهم{شرطی‌سازی سطح}: یک یا چند مقدار عددی که سطح روشنایی هدف در اندام‌های شاخص را مشخص می‌کند.
\end{itemize}

شرطی‌سازی سطح نقش ویژه‌ای دارد و باید به‌دقت تفسیر شود. این مقادیر از تصویر کنتراست‌دار \emph{واقعی} خوانده می‌شوند و بنابراین در زمان استقرار بالینی در دسترس نیستند. هدف از این پیکربندی، ارائه‌ی یک روش کاربردی نیست؛ هدف، اندازه‌گیری یک \emph{سقف} است: اگر مدل هنگام دریافت سطح واقعی بتواند آن را بازتولید کند، آن‌گاه ناتوانی‌اش در حالت عادی ناشی از کمبود \emph{اطلاعات ورودی} است و نه کمبود ظرفیت. این تمایز، محور استدلال فصل~\ref{chap:experiments} است.

\subsection*{استنتاج در سطح حجم}

استنتاج به‌صورت کاشی‌کاری با هم‌پوشانی و آمیختن\footnote{\lr{Blending}} وزنی انجام می‌شود تا مرز وصله‌ها در خروجی دیده نشود. خروجی به \lr{HU} بازگردانده و روی شبکه‌ی اصلی هر مورد آزمون به‌صورت \lr{NIfTI} نوشته می‌شود. این قرارداد خروجی برای همه‌ی مدل‌ها—از جمله مدل‌های بازپیاده‌سازی‌شده از ادبیات—یکسان است، و همین یکسانی است که امکان امتیازدهی همه‌ی آن‌ها با یک ابزار واحد را فراهم می‌کند.

\subsection*{مدل‌سازی عدم‌قطعیت: انتشار و رگرسیون ناهم‌واریانس}
\label{sec:m-uncertainty}

یک رگرسیون‌گر آموزش‌دیده با زیان $L_1$، در بهترین حالت \emph{میانه‌ی شرطی} $\mathbb{E}[y \mid x]$ را پیش‌بینی می‌کند. اگر رابطه‌ی $x \mapsto y$ قطعی نباشد—یعنی اگر $y$ به متغیرهایی وابسته باشد که در $x$ مشاهده‌پذیر نیستند—آن‌گاه خطای باقی‌مانده تقلیل‌ناپذیر است و هیچ تغییری در معماری یا تابع زیان آن را کاهش نمی‌دهد.

در مسئله‌ی حاضر، نامزد چنین متغیرهای مشاهده‌نشده‌ای روشن است: دوز تزریق، زمان‌بندی بولوس و برون‌ده قلبی، که هیچ‌یک در اسکن بدون‌کنتراست قابل‌مشاهده نیستند. به بیان صوری،
\[
y = f(x, Z), \qquad Z = \text{(دوز، زمان‌بندی، فیزیولوژی)} \ \text{مشاهده‌نشده} .
\]
اگر این صورت‌بندی درست باشد، مدل قطعی نه تنها خطا دارد، بلکه \emph{باید} داشته باشد؛ و آن‌چه باید سنجیده شود نه کاهش خطا، بلکه \emph{کالیبراسیون} توزیع پیش‌بینی‌شده است.

دو خانواده‌ی مدل برای آزمودن این ایده به کار می‌روند:

\paragraph{رگرسیون ناهم‌واریانس.}
سر خروجی مولد به دو کانال گسترش می‌یابد که میانگین $\mu$ و لگاریتم واریانس $\log \sigma^2$ را در هر وکسل پیش‌بینی می‌کنند، و آموزش با درست‌نمایی منفی گاوسی انجام می‌شود~\cite{kendall2017uncertainties}. این مدل با هزینه‌ی محاسباتی ناچیز، یک تخمین عدم‌قطعیت کالیبره‌شدنی می‌دهد و بنابراین \emph{مبنایی} است که مدل انتشار باید از آن بهتر عمل کند تا هزینه‌اش توجیه شود.

\paragraph{مدل انتشار.}
مدل انتشار~\cite{ho2020denoising} به‌جای $\mathbb{E}[y \mid x]$، خودِ توزیع $p(y \mid x)$ را مدل می‌کند و از این نظر ابزار نظری درست برای یک باقی‌مانده‌ی تصادفی است. پیاده‌سازی این پایان‌نامه از پارامترسازی $v$~\cite{salimans2022progressive}، نمونه‌بردار \lr{DDIM}~\cite{song2021ddim}، هدایت بدون‌طبقه‌بند~\cite{ho2022classifier}، وزن‌دهی زیان بر پایه‌ی نسبت سیگنال به نویز کمینه~\cite{hang2023minsnr} و میانگین متحرک نمایی وزن‌ها بهره می‌گیرد.

نکته‌ی حیاتی—که پیش از اجرای آزمایش‌ها ثبت شده است—آن است که یک نمونه‌ی صادق از توزیع شرطی، بر حسب یک معیار نقطه‌ای، \emph{بدتر} از میانگین شرطی امتیاز می‌گیرد. برای یک توزیع گاوسی با انحراف معیار $s$، پیش‌بین میانگین خطای مطلق $0.798\,s$ و یک نمونه‌ی منفرد خطای $1.128\,s$ دارد؛ یعنی نسبتی حدود $\sqrt{2}$. این افت، نقص مدل انتشار نیست بلکه بهای صادقانه‌ی مدل‌کردن عدم‌قطعیت است، آن‌گاه که با معیاری نقطه‌ای سنجیده شود. ثبت این انتظار \emph{پیش از} آزمایش، شرط لازم برای آن است که نتیجه قابل ابطال باشد.

\subsection*{مدل‌های مبنا از ادبیات}

برای آن‌که ادعاهای این پایان‌نامه در برابر وضعیت موجود سنجیده شوند، مجموعه‌ای از مدل‌های منتشرشده روی \emph{همان} تقسیم داده، \emph{همان} موارد آزمون و با \emph{همان} ابزار امتیازدهی بازآموزی و ارزیابی می‌شوند:

\begin{itemize}
\item \lr{pix2pix} و \lr{pix2pixHD}~\cite{isola2017image,wang2018pix2pixhd} — مبنای ترجمه‌ی جفت‌شده.
\item \lr{CycleGAN}~\cite{zhu2017unpaired} — شاهد \emph{جفت‌نشده}. داده‌ی ما جفت‌شده است، اما این مدل عمداً در حالت جفت‌نشده آموزش می‌بیند تا هزینه‌ی چشم‌پوشی از جفت‌شدگی اندازه‌گیری شود.
\item \lr{ResViT}~\cite{dalmaz2022resvit} و \lr{CyTran}~\cite{Ris21} — مبناهای مبتنی‌بر ترانسفورمر.
\item \lr{SwinUNETR}~\cite{hatamizadeh2021swin} و \lr{TransUNet}~\cite{chen2021transunet} — این دو \emph{معماری}‌اند و نه پایپ‌لاین سنتز؛ هر دو در اصل برای بخش‌بندی عرضه شده‌اند. شبکه‌ی آن‌ها بدون تغییر نگه داشته و حلقه‌ی آموزش با هدف یکسان با \lr{ResViT} فراهم می‌شود تا مقایسه، \emph{معماری} را جدا کند و نه بهینه‌سازی را.
\item \lr{Ea-GAN}~\cite{yu2019eagan} — شبکه‌ی مولد-متخاصم آگاه از لبه.
\item \lr{SynDiff}~\cite{ozbey2023syndiff} — مدل انتشار خصمانه. این مدل به‌دلیل هزینه‌ی محاسباتی از اجرای نهایی کنار گذاشته شد؛ دلیل و شواهد زمان‌سنجی در فصل~\ref{chap:experiments} گزارش می‌شود.
\end{itemize}

باید تصریح شود که بازآموزی این مدل‌ها روی داده و مقیاس این پایان‌نامه، اعداد مقاله‌های اصلی آن‌ها را بازتولید \emph{نمی‌کند}؛ آن‌چه گزارش می‌شود یک مقایسه‌ی کنترل‌شده بر داده‌ی یکسان است، نه یک بازتولید.

\section{پروتکل ارزیابی}
\label{sec:m-eval}

\subsection*{چرا معیارهای متعارف کافی نیستند}

معیارهای رایج در ادبیات سنتز تصویر پزشکی—\lr{PSNR}، \lr{SSIM} و خطای مطلق میانگین—در این مسئله گمراه‌کننده‌اند، و این نکته یکی از یافته‌های روش‌شناختی این پایان‌نامه است.

دلیل نخست ساختاری است: پنجره‌ی \lr{HU} انتخاب‌شده، هوا، ریه و چربی را به صفر و استخوان را به یک اشباع می‌کند، و این نواحی بخش بزرگی از حجم را تشکیل می‌دهند. هر معیاری که روی کل تصویر میانگین بگیرد، عمدتاً روی جرمی میانگین می‌گیرد که همه‌ی مدل‌ها در آن بی‌خطا هستند؛ حتی یک کپی همانی از تصویر بدون‌کنتراست بخش عمده‌ی امتیاز بهترین مدل را کسب می‌کند.

دلیل دوم عمیق‌تر است: این معیارها با میانگین شرطی کمینه می‌شوند. بنابراین یک خروجی \emph{تار} را بالاتر از یک نمونه‌ی دارای بافت \emph{صحیح} امتیاز می‌دهند. مرتب‌سازی مدل‌ها بر اساس این معیارها، آن‌ها را بر اساس میزان بافتی که تلاش می‌کنند تولید کنند مرتب می‌کند—در جهت اشتباه.

بنابراین این معیارها گزارش می‌شوند اما \emph{ثانویه} تلقی می‌شوند، و ادعاها بر معیارهای زیر بنا می‌شود.

\subsection*{معیار اصلی: وفاداری \lr{HU} در سطح اندام}

معیار سرصفحه‌ی این پایان‌نامه، میانگین قدرمطلق خطای میانه‌ی \lr{HU} در سطح اندام است که روی مجموعه‌ای از اندام‌های شکمی محاسبه می‌شود. نکته‌ی مهم آن است که این معیار با همان استخراج ویژگی سطح‌اندامِ مدل تشخیص فاز (بخش~\ref{sec:m-phase}) محاسبه می‌شود. به این ترتیب مؤلفه‌ی نخست روش، ابزار سنجش مؤلفه‌ی سوم می‌شود، و معیار به کمیتی متصل می‌ماند که معنای بالینی دارد: میزان روشنایی هر اندام، همان چیزی که فاز کنتراست را تعریف می‌کند.

در کنار آن، خطای مطلق و \lr{PSNR} و \lr{SSIM} \emph{محدودشده به ناحیه‌ی اندام} گزارش می‌شود که مشکل نخستِ بالا را برطرف می‌کند.

\subsection*{معیار بازیابی سطح: تشخیص مولد از میانگین‌گیر}

دو کمیت برای تشخیص یک مولد واقعی از یک میانگین‌گیر تعریف می‌شود. میانه‌ی \lr{HU} هر اندام در تصویر تولیدشده، بر همان کمیت در تصویر واقعی، در سراسر موارد آزمون رگرس می‌شود:

\begin{itemize}
\item $\beta$، شیب این رگرسیون. مقدار $\beta = 1$ یعنی مدل سطح واقعی هر مورد را دنبال می‌کند؛ مقدار $\beta = 0$ یعنی مدل صرف‌نظر از ورودی، میانگین جمعیت را تولید می‌کند.
\item نسبت واریانس $\mathrm{var}(\text{تولیدشده}) / \mathrm{var}(\text{واقعی})$، با هدف $1$. مقدار نزدیک صفر یعنی خروجی کم‌پراکنده\footnote{\lr{Under-dispersed}} است.
\end{itemize}

این دو کمیت دقیقاً همان چیزی را می‌سنجند که معیارهای نقطه‌ای پنهان می‌کنند: یک مدل می‌تواند خطای \lr{HU} قابل‌قبولی داشته باشد و در عین حال نسبت واریانسش نزدیک صفر باشد، یعنی هرگز سطح کنتراست هیچ بیماری را واقعاً پیش‌بینی نکرده باشد.

\subsection*{معیارهای بافت و پیوستگی}

چهار کمیت نسبتی برای سنجش بافت و پیوستگی به کار می‌رود: نسبت انرژی طیفی بسامد بالا، فاصله‌ی واسرشتاین میان توزیع بزرگی گرادیان تولیدشده و واقعی، انرژی درز در مرز کاشی‌ها، و سوسوزنی\footnote{\lr{Flicker}} میان‌برشی.

نکته‌ی تفسیری مهم آن است که این‌ها \emph{نسبت}‌اند و هدفشان مقدار $1$ است، نه بیشینه. مقدار کمتر از یک به معنای تاری و مقدار بیشتر از یک به معنای نویز یا بافت توهمی\footnote{\lr{Hallucinated}} است؛ هر دو شکست‌اند و باید بر حسب فاصله از یک رتبه‌بندی شوند.

معیارهای ادراکی \lr{LPIPS}~\cite{zhang2018lpips} و \lr{FID}~\cite{heusel2017fid} نیز صرفاً برای \emph{قابلیت مقایسه با ادبیات} گزارش می‌شوند و مبنای هیچ ادعایی نیستند: هر دو شبکه‌های از پیش آموزش‌دیده روی تصاویر طبیعی را بر تصویر سی‌تی خاکستری اجرا می‌کنند و اعتبار ترتیبشان در این دامنه تأیید نشده است.

\subsection*{معیارهای کالیبراسیون}

برای مدل‌های احتمالاتی بخش~\ref{sec:m-uncertainty}، دو کمیت افزوده می‌شود: امتیاز احتمال رتبه‌ای پیوسته\footnote{\lr{Continuous ranked probability score, CRPS}}~\cite{gneiting2007scoring} و پوشش بازه‌ی پیش‌بینی ۹۰ درصد، که مقدار هدفش $0.90$ است. انحراف از این مقدار، مستقیماً بدکالیبراسیون را اندازه می‌گیرد.

\subsection*{ارزیابی پایین‌دستی: شکاف بخش‌بندی}
\label{sec:m-downstream}

هیچ‌یک از معیارهای بالا، پرسش اصلی این پایان‌نامه را پاسخ نمی‌دهد. آن پرسش یک ارزیابی جداگانه می‌طلبد که به‌صورت زیر تعریف می‌شود.

مرجع، خروجی بخش‌بندی خودکار روی تصویر کنتراست‌دار \emph{واقعی} است. سه بازو در برابر این مرجع سنجیده می‌شوند:

\begin{itemize}
\item بخش‌بندی روی تصویر بدون‌کنتراست واقعی — این \emph{خودِ شکاف} است.
\item بخش‌بندی روی تصویر کنتراست‌دار \emph{تولیدشده} توسط هر مدل.
\item بخش‌بندی روی تصویر تولیدشده توسط مدلِ شرطی‌شده با سطح واقعی — که \emph{سقف} قابل‌دستیابی را می‌دهد.
\end{itemize}

معیارها همان ضریب دایس فرسوده و فاصله‌ی هاسدورف صدک ۹۵ بخش~\ref{sec:m-reg} هستند، تا فصل ثبت تصویر و فصل سنتز یک تعریف واحد از دایس داشته باشند. اندام‌ها به همان چهار گروه دسته‌بندی می‌شوند و گروه \emph{اسکلتی} بار دیگر نقش شاهد منفی را دارد: استخوان با کنتراست تغییر نمی‌کند، پس اگر بازوی بدون‌کنتراست روی استخوان نیز افت دایس نشان دهد، آن‌چه اندازه گرفته شده تفاوت کیفیت کلی تصویر است و نه اثر کنتراست، و بقیه‌ی جدول را نمی‌توان به کنتراست نسبت داد.

محدودیت این طرح باید صریح بماند: مرجع، حاشیه‌نویسی انسانی نیست بلکه خروجی خودکار روی تصویر واقعی است. بنابراین یک تصویر تولیدشده می‌تواند با شبیه‌شدن به یک تصویر کنتراست‌دار \emph{از دید همان ابزار بخش‌بندی} امتیاز بگیرد. این محدودیت هرجا که این اعداد ظاهر شوند تکرار خواهد شد؛ شاهد منفی اسکلتی و کف بازوی بدون‌کنتراست، آن چیزی است که سیگنال باقی‌مانده را قابل‌تفسیر نگه می‌دارد.

\subsection*{انضباط آماری}

سه قاعده در سراسر فصل~\ref{chap:experiments} اعمال می‌شود و ذکرشان در این‌جا ضروری است، چون بدون آن‌ها بخشی از نتایج گزارش‌شده متفاوت می‌بود.

\begin{itemize}
\item \مهم{کف نویز بذر.} هر پیکربندی کلیدی با سه بذر تصادفی اجرا و انحراف معیار میان بذرها محاسبه می‌شود. تفاوتی که از دروازه‌ی $2\sigma$ کوچک‌تر باشد نتیجه به شمار نمی‌آید. این قاعده در جریان همین پژوهش یک ادعای سرصفحه را باطل کرد—ادعایی که در یک بذر معنادار به نظر می‌رسید و در سه بذر نه‌تنها ناپدید شد بلکه علامتش برگشت.
\item \مهم{آزمون‌های زوجی در سطح مورد.} پراکندگی میان بیماران بر تفاوت میان مدل‌ها غلبه دارد. بنابراین مقایسه‌ها از طریق آزمون زوجی روی موارد مشترک انجام می‌شود و نه از طریق مقایسه‌ی میانگین‌ها و انحراف معیارها.
\item \مهم{تفاضل پیکربندی.} هر اجرا باید با اجرای پیشین خود تنها در کلیدهایی تفاوت داشته باشد که عمداً تغییر داده شده‌اند. این ساده‌ترین محافظ در برابر مخدوش‌کننده‌هاست.
\end{itemize}

\section{جمع‌بندی}
\label{sec:m-summary}

این فصل سه مؤلفه‌ی روش پیشنهادی را ارائه کرد. مدل تشخیص فاز، مسئله را از حجم خام به بردار کوچکی از میانه‌های \lr{HU} در سطح اندام تقلیل می‌دهد و در ادامه، علاوه بر نقش دروازه‌بانی داده، ابزار اندازه‌گیری ارزیابی سنتز نیز می‌شود. پایپ‌لاین ثبت تصویر چندگذره، به‌جای وزن‌دهی به نواحی پایدار، \emph{دامنه‌ی} محاسبه‌ی معیار تراز را به آن‌ها محدود می‌کند و با مقداردهی اولیه‌ی بسته بر پایه‌ی نشانه‌های استخوانی از کمینه‌های محلی پرهیز می‌کند. مدل‌های سنتز، از مولد قطعی تا مدل انتشار و رگرسیون ناهم‌واریانس، همگی روی جفت‌های حاصل از این ثبت تصویر و روی یک تقسیم منجمد آموزش می‌بینند.

پروتکل ارزیابی، خود بخشی از سهم این پایان‌نامه است: نشان داده شد که چرا معیارهای متعارف سطح‌پیکسل در این مسئله قادر به رتبه‌بندی مدل‌ها نیستند، و مجموعه‌ای از معیارهای سطح‌اندام، بازیابی سطح، بافت و کالیبراسیون به‌جای آن‌ها تعریف شد. سرانجام، ارزیابی پایین‌دستی بخش‌بندی تعریف شد که تنها سنجه‌ای است که به پرسش عنوان این پایان‌نامه پاسخ می‌دهد.

فصل بعد نتایج تجربی هر یک از این مؤلفه‌ها را گزارش می‌کند.
